Soit $\K$ un sous-corps de $\C$ et $E$ un $\K$-espace vectoriel.
 
\section{Sous-espaces stables}

\begin{definition}{Endomorphisme induit}{}
    \index{Endomorphisme!induit}
    Soit $u \in \mathscr L(E)$ et $F$ un sous-espace vectoriel de $E$. On dit que 
    $F$ est stable par $u$ ssi $\forall x \in F, u(x) \in F$.

    En ce cas on peut définir $\fonction{\tilde u}{F}{F}{x}{u(x)}$ qui est une application 
    linéaire, appelée endomorphisme induit par $u$ sur $F$.
\end{definition}

\subsection{Cas de la dimension finie}

\begin{proposition}{Supplémentaire et base adaptée}{}
    \index{Supplementaire@Supplémentaire}
    \index{Base adaptée}
    Soit $E$ un $\K$-espace vectoriel de dimension finie $n$, $F$ un sous-espace vectoriel de $E$
    de dimension $p$, avec $1 \leqslant p \leqslant n - 1$. 

    Soit $B_1$ une base de $F$, et complétons la en $B = B_1 \sqcup B_2$ une base de $E$. Alors 
    $\Vect (B_2) = G$ est un supplémentaire de $F$ dans $E$ et 
    $B$ est adaptée à la décomposition $E = F \oplus G$.

\end{proposition}

\begin{proposition}{Caractérisation matricielle de la stabilité}{}
    En reprenant les notations précédentes, soit $u \in \mathscr L(E)$ et $A = M_B(u) \in \mathscr M_n(\K)$. 
    Alors $F$ est stable par $u$ ssi 
    $\exists A' \in \mathscr M_{p}(\K), C \in \mathscr M_{p, n - p}(\K)$ et $D \in \mathscr M_{n - p}(\K)$ tels que
    \[
        A = \begin{pmatrix}
            A' & C \\
            0 & D
        \end{pmatrix}
    \]

    En ce cas en notant $\tilde u$ l'endomorphisme induit par $u$ sur $F$, on a $A' = M_{B_1}(\tilde u)$.
\end{proposition}

\begin{proof}
    Supposons que $F$ est stable par $u$. Alors pour tout $j \in \ninter 1 p $ en notant 
    $B_1 = (e_1, \ldots, e_p)$ on a $e_i \in F$ donc $u(e_j) \in F$. 

    Donc les coordonnées de $u(e_j)$ selon les $e_i$ pour $i > p$ sont toutes nulles, donc 
    $A$ a la forme voulue.

    \lvspace Réciproquement, si $A$ a la forme voulue, pour tout $j \in \ninter 1 p$ on a
    $u(e_j) \in F$ donc pour tout $x = \sum_{i = 1}^p x_i e_i \in F$ on a
    \[
        u(x) = \sum_{i = 1}^p x_i u(e_i) \in F
    \]
    donc $F$ est stable par $u$ et donc les colonnes de $A'$ sont bien les coordonnées dans $B_1$ 
    des vecteurs de $u(B_1)$.
\end{proof}

% \restaurergeometrie

\section{\'Eléments propres d'un endomorphisme ou d'une matrice carrée }
\subsection{Vecteur propre, valeur propre, sous-espace propre }

Soit $E$ un $\K$-espace vectoriel non nul et $u \in \mathscr L(E)$.

\begin{definition}{Valeur propre}{}
    \index{Valeur propre!d'un endomorphisme}
    Soit $\lambda \in \K$. On dit que $\lambda$ est une valeur propre de $u$ ssi
    $\exists x \in E$ tel que $\begin{cases} x \neq 0 \\ u(x) = \lambda x \end{cases}$.
\end{definition}

\begin{definition}{Vecteur propre}{}
    \index{Vecteur propre!d'un endomorphisme}
    Soit $x \in E$. On dit que $x$ est un vecteur propre de $u$ ssi
    $\exists \lambda \in \K$ tel que $\begin{cases} x \neq 0 \\ u(x) = \lambda x \end{cases}$.
\end{definition}

\begin{remarque}
    Soient $\lambda \in \K$ et $x \in E$. Si $x \neq 0$ et $u(x) = \lambda x$, alors 
    on dit que $\lambda$ est \emph{la} valeur propre associée au vecteur propre $x$ et 
    que $x$ est \emph{un} vecteur propre associé à la valeur propre $\lambda$.

    Si $x$ est un vecteur propre pour la valeur propre $\lambda$, il n'est pas unique 
    car $2x$ l'est aussi par exemple.
\end{remarque}

\begin{definition}{Espace propre}{}
    \index{Espace propre!d'un endomorphisme}
    Soit $\lambda \in \K$ une valeur propre de $u$. On appelle \emph{espace propre} 
    associé l'espace vectoriel 
    \begin{align} E_\lambda(u) & = \{ x \in E \mid u(x) = \lambda x \} \\
                        & = \Ker (u - \lambda \Id_E) 
    \end{align}
    Il est donc formé des vecteurs propres de $u$ associés à la valeur propre $\lambda$,
    ainsi que du vecteur nul.
\end{definition}

\begin{remarque}
    Si $\lambda$ n'est pas une valeur propre de $u$, alors $E_\lambda(u) = \{0\}$.
\end{remarque}

\begin{definition}{Spectre}{}
    \index{Spectre!d'un endomorphisme}
    Si $E$ est de dimension finie, l'ensemble des valeurs propres de $u$ est appelé spectre de $u$ 
    et noté $\Sp(u)$.
\end{definition}

\begin{definition}{Équation aux éléments propres}{}
    \index{Equation aux elements propres@\'Equation aux éléments propres!d'un endomorphisme}
    Soit $u \in \mathscr L(E)$. On appelle \emph{équation aux éléments propres} de $u$
    l'équation \\ 
    $u(x) = \lambda x$ où $x \in E \setminus \{0\}$ et $\lambda \in \K$ sont 
    les inconnues. $(\lambda, x)$ est une solution de cette équation ssi $\lambda$ est une valeur propre 
    de $u$ et $x$ un vecteur propre associé.
\end{definition}

\subsection{Exemples}
\subsubsection{Homothéties}

\begin{exemple}
    Soit $\lambda \in \K$ et $u = \lambda \Id_E$ l'homothétie de rapport $\lambda$. 
    Alors tout vecteur de $E$ est vecteur propre pour la valeur propre $\lambda$.
    Donc $E_\lambda(u) = E$.

    \avspace Réciproquement, si tout vecteur non nul de $E$ est propre pour un endomorphisme $u$, 
    alors $u$ est une homothétie (fais dans un exo du TD du chapitre 11).
\end{exemple}

\subsubsection{Projections et symétries}

\begin{exemple}
    soient $F$ et $G$ deux sous-espaces vectoriels de $E$ tels que $E = F \oplus G$.
    Soit $\pi$ la projection sur $F$ parallèlement à $G$ et $s = 2 \pi - \Id_E$ la symétrie associée. 

    Ainsi si $x = x_F + x_G \in E$, alors $\pi(x) = x_F$ et $s(x) = x_F - x_G$.

    On suppose que $F$ et $G$ sont non nuls. Alors les valeurs propres de $\pi$ sont $0$ et $1$ 
    et les espaces propres associés sont $E_0(\pi) = G$ et $E_1(\pi) = F$.

    Les valeurs propres de $s$ sont $-1$ et $1$ et les espaces propres associés sont 
    $E_{-1}(s) = G$ et $E_1(s) = F$.

    \uindent{Preuve }

    \begin{proof} 
        Pour trouver vecteurs et valeurs propres on écrit l'équation aux éléments propres.

        Soit $\lambda \in \K, x \in E \setminus \{0\}$ tels que $\pi(x) = \lambda x$.
        En écrivant $x = x_F + x_G$ avec $x_F \in F$ et $x_G \in G$, on a
        
        $\pi(x) = \lambda x $ ssi $x_F = \lambda (x_F + x_G)$ ssi $\begin{cases}
        x_F = \lambda x_F \\
        0 = \lambda x_G
        \end{cases}$ ssi $\begin{cases}
        (1 - \lambda) x_F = 0 \\
        \lambda x_G = 0
        \end{cases}$ ssi $\begin{cases}
            \lambda = 1 \text{ ou } x_F = 0 \\
            \lambda = 0 \text{ ou } x_G = 0
        \end{cases}$.

        Donc $\pi(x) = \lambda x$ ssi $\begin{cases}
            \lambda = 1 &\text{ et } x_G = 0 \\
            & \text{ ou } \\
            \lambda = 0 &\text{ et } x_F = 0
        \end{cases}$

        \avspace De même, $s(x) = \lambda x$ ssi $(2 \pi - \Id_E)(x) = \lambda x$ ssi
        $2 \pi(x) = (\lambda + 1) x$ ssi $\pi(x) = \frac{\lambda + 1}{2} x$ 
        ssi $\begin{cases}
            \frac{\lambda + 1}{2} = 1 &\text{ et } x_G = 0 \\
            & \text{ ou } \\
            \frac{\lambda + 1}{2} = 0 &\text{ et } x_F = 0
        \end{cases}$ ssi $\begin{cases}
            \lambda = 1 &\text{ et } x_G = 0 \\
            & \text{ ou } \\
            \lambda = -1 &\text{ et } x_F = 0
        \end{cases}$.
    \end{proof}
\end{exemple}

\subsubsection{Autres exemples }

\marginenote{Important !}{Ce point peut s'avérer utile en exercice.}

\avspace

\begin{exemple}
    Soit $u \in \mathscr L(E)$. Alors $0$ est valeur propre de $u$ ssi $u$ n'est pas injectif.
    En effet, $0$ est valeur propre ssi $\exists x \in E \setminus \{0\}$ tel que $u(x) = 0$ 
    ssi $\Ker(u) \neq \{0\}$ ssi $u$ n'est pas injectif.

    Dans ce cas, l'espace propre associé est $E_0(u) = \Ker(u)$.
\end{exemple}

Un endomorphisme peut ne posséder aucune valeur propre.

\begin{exemple}
    Par exemple, $u = $ rotation d'angle $\frac \pi 2$ dans le plan euclidien orienté. 
    $u$ est alors associé dans la base canonique à la matrice
    \[
        J = \begin{pmatrix}
            0 & -1 \\
            1 & 0
        \end{pmatrix}
    \]

    en ce cas pour tout $\lambda \in \R$, $\det(u - \lambda \Id_E) = \det \begin{pmatrix}
        -\lambda & -1 \\
        1 & -\lambda
    \end{pmatrix} = \lambda^2 + 1 \neq 0$ donc $u - \lambda \Id_E$ est bijectif, donc 
    $\Ker(u - \lambda \Id_E) = \{0\}$ donc pas de vecteur propre. 
\end{exemple}

Un endomorphisme peut avoir une infinité de valeurs propres :

\marginenote{Détail}{Cela ne fonctionne que parce qu'il s'agit d'un espace vectoriel de dimension infinie.}

\begin{exemple}
    Posons $E = \mathscr C^\infty(\R, \R)$ l'espace vectoriel des fonctions 
    infiniment dérivables de $\R$ dans $\R$ et $\fonction u E E f {f'}$ la dérivation. 

    $\forall \lambda \in \R$, $\lambda$ est valeur propre pour les vecteurs propres 
    $$\fonction{f_\lambda}{\R}{\R}{x}{e^{\lambda x}}$$ car
    $u(f_\lambda) = f_\lambda' = \lambda f_\lambda$.
\end{exemple}

\begin{exemple}
    $E = \R^{\N}$ et $\fonction s E E {(a_n)_{n \in \N}}{(a_{n + 1})_{n \in \N}}$ l'application de décalage à gauche.

    Tout $\lambda \in \R$ est valeur propre pour le vecteur propre du terme général $a_n = \lambda^n$.
\end{exemple}

\subsection{Droites stables }

\begin{proposition}{Droite stable et vecteur propre}{}
    Soit $u \in \mathscr L(E)$, $a \in E \setminus \{0\}$ et $D = \Vect(a)$ une droite de $E$. 

    Alors $D$ est stable par $u$ ssi $a$ est vecteur propre de $u$.

    Ainsi les vecteurs non nuls d'une droite stable sont tous des vecteurs propres. 
\end{proposition}

\begin{proof}
    $\Rightarrow$ Supposons que $D$ est stable par $u$. Alors $u(a) \in D$ donc 
    $\exists \lambda \in \K$ tel que $u(a) = \lambda a$, donc $a$ est vecteur propre de $u$.
    Or $a \neq 0$ donc $a$ est bien un vecteur propre de $u$.

    \lvspace $\Leftarrow$ Réciproquement, si $a$ est vecteur propre de $u$, alors 
    $\exists \lambda \in \K$ tel que $u(a) = \lambda a$. 

    Donc pour tout $x \in D$, $\exists \mu \in \K$ tel que $x = \mu a$ et donc 
    $u(x) = u(\mu a) = \mu u(a) = \mu \lambda a \in D$. Donc $D$ est stable par $u$.
\end{proof}

\subsection{Lien entre sous-espaces propres }

\begin{theorement}{Nature famille de vecteurs propres associée à des valeurs propres distinctes}{}
    Toute famille de vecteurs propres pour un même endomorphisme 
    associés à des valeurs propres deux à deux distinctes est libre. 
\end{theorement}

\idee{
    Par récurrence sur le cardinal de la famille.
}

\begin{proof}
    On peut ramener la preuve au cas d'une famille finie car une famille est libre 
    ssi toute sous famille finie est libre. 

    Raisonnons par récurrence sur $k$ le cardinal de la famille. 

    \uindent{Initialisation } $k = 1$

    \svspace soit $x$ un vecteur propre pour $u$ alors $x \neq 0$ donc la famille $(x)$ est libre.

    \uindent{Hérédité } Supposons le résultat pour toute famille de $k$ vecteurs. 

    \svspace Soient $x_1, \ldots, x_{k + 1} \in E$ des vecteurs propres de $u$ pour les 
    valeurs propres $\lambda_1, \ldots, \lambda_{k + 1}$.

    Soient $\alpha_1, \ldots, \alpha_{k + 1} \in \K$ et supposons que 
    $\displaystyle \sum_{i = 1}^{k + 1} \alpha_i x_i = 0$.

    Appliquons $u$ à cette relation :
    \[
        \sum_{i = 1}^{k + 1} \alpha_i u(x_i) = \sum_{i = 1}^{k + 1} \alpha_i \lambda_i x_i = 0
    \]
    et alors en combinant ces deux relations (en prenant $- \lambda_{k+1}$ fois la première), on obtient :
    \[
        \sum_{i = 1}^{k} \alpha_i (\lambda_i - \lambda_{k + 1}) x_i = 0
    \]

    Or par hypothèse de récurrence, la famille $(x_1, \ldots, x_k)$ est libre et 
    les $\lambda_i$ sont deux à deux distincts donc $\forall i \in \ninter 1 k, \lambda_i - \lambda_{k + 1} \neq 0$.

    Donc $\forall i \in \ninter 1 k, \alpha_i = 0$ et donc $\alpha_{k + 1} = 0$. 
    Ainsi la famille $(x_1, \ldots, x_{k + 1})$ est libre.
\end{proof}

\begin{exemple}
    Pour $\alpha \in \C$ notons $\fonction{f_\alpha}{\R}{\C}{x}{e^{\alpha x}}$.
\end{exemple}

\begin{theoreme}{Somme de sous-espaces propres}{}
    Une somme de sous-espaces propres associés à des valeurs propres deux à deux distinctes 
    d'un même endomorphisme est directe. 
\end{theoreme}

\idee{
    Ne pas oublier que les espaces propres contiennent le vecteur nul.
}

\begin{proof}
    Soit $u \in \mathscr L(E)$ et $\lambda_1, \ldots, \lambda_k$ des valeurs propres deux à deux distinctes de $u$
    et $S = E_{\lambda_1}(u) + \cdots + E_{\lambda_k}(u)$.

    Soient $x_1 \in E_{\lambda_1}(u), \ldots, x_k \in E_{\lambda_k}(u)$ tels que
    $\displaystyle \sum_{i = 1}^k x_i = 0$.

    Supposons que certains des $x_i$ sont non nuls, par exemple (quitte à renuméroter les indices)
    $x_1, \dots, x_p$ où $p \leqslant k$ alors $\displaystyle \sum_{i = 1}^p x_i = 0$.

    Donc la famille $(x_1, \ldots, x_p)$ est liée, ce qui contredit le théorème précédent. 

    Donc $\forall i \in \ninter 1 k, x_i = 0$ et la somme est directe.
\end{proof}

\begin{theoreme}{Nombre de valeurs propres d'un endomorphisme}{}
    Soit $E$ un $\K$-espace vectoriel de dimension finie $n$ et soit $u \in \mathscr L(E)$.

    Alors $u$ possède au plus $n$ valeurs propres c'est-à-dire que $\Sp(u)$ est fini et 
    $\card(\Sp(u)) \leqslant n$.
\end{theoreme}

\idee{
    Dimension.
}

\begin{proof}
    Supposons par l'absurde que $u$ possède $n + 1$ valeurs propres deux à deux distinctes
    notées $\lambda_1, \ldots, \lambda_{n + 1}$.

    Alors $S = E_{\lambda_1}(u) + \cdots + E_{\lambda_{n + 1}}(u)$ est une somme directe
    donc 
    
    $\displaystyle \dim(S) = \sum_{i = 1}^{n + 1} \dim(E_{\lambda_i}(u)) \geqslant n + 1$ car
    chaque espace propre contient au moins un vecteur non nul.

    Absurde car $S$ est un sous-espace de $E$ donc $\dim(S) \leqslant \dim(E) = n$.
\end{proof}

\subsection{Stabilité de sous-espaces }

\begin{proposition}{Commutation et stabilité des sous-espaces propres}{}
    Soient $u, v \in \mathscr L(E)$ tel que $u \circ v = v \circ u$.

    Alors tout sous-espace propre de $u$ est stable par $v$.
\end{proposition}

\begin{proof}
    Soit $F$ un sous-espace propre de $u$. 

    Soit $\lambda \in \K$ valeur propre de $u$ telle que $F = E_\lambda(u)$.

    Soit $x \in F$, montrons que $v(x) \in F$.

    \begin{align}
        u(v(x)) & = (u \circ v)(x) \\
                & = (v \circ u)(x) \\
                & = v(u(x)) \\
                & = v(\lambda x) \\
                & = \lambda v(x)
    \end{align}

    Donc $v(x) \in F$.
\end{proof}

\begin{proposition}{Commutation et stabilité du noyau et de l'image}{}
    Soient $u, v \in \mathscr L(E)$ tel que $u \circ v = v \circ u$. Alors $\Ker(u)$ et
    $\im(u)$ sont stables par $v$.
\end{proposition}

\begin{proof}
    $\Ker u = E_0(u)$ donc stable par $v$ d'après la proposition précédente.

    Soit $y \in \im(u)$ alors $\exists x \in E$ tel que $y = u(x)$. Alors 
    $v(y) = v(u(x)) = u(v(x)) \in \im(u)$ donc $\im(u)$ est stable par $v$.
\end{proof}

\begin{proposition}{Cas des sous-espaces propres}{}
    Soit $u \in \mathscr L(E)$ et $\lambda \in \K$ une valeur propre de $u$. 
    Alors l'espace propre $E_\lambda(u)$ est stable par $u$ et l'endomorphisme induit par $u$ 
    sur $E_\lambda(u)$ est l'homothétie de rapport $\lambda$.
\end{proposition}

\begin{proof}
    Immédiate par la proposition précédente. 
\end{proof}

\subsection{Eléments propres d'une matrice carrée }

\begin{definition}{}{}
    Soit $A \in \mathscr M_n(\K)$. On définit l'endomorphisme de $\mathscr M_{n, 1}(\K)$
    identifié à $\K^n$ canoniquement associé à $A$ par 
    $$\fonction{u_A}{\K^n = \mathscr M_{n, 1}(\K)}{\K^n}{X}{AX}$$
\end{definition}

\begin{definition}{Valeur propre d'une matrice}{}
    \index{Valeur propre!d'une matrice}
    Soit $\lambda \in \K$. On dit que $\lambda$ est une valeur propre de $A$ ssi
    c'est une valeur propre de l'endomorphisme $u_A$, c'est à dire s'il existe 
    $X \in \K^n \setminus \{0\}$ tel que $AX = \lambda X$.
\end{definition}

\begin{definition}{Vecteur propre d'une matrice}{}
    \index{Vecteur propre!d'une matrice}
    Soit $X \in \K^n$. On dit que $X$ est un vecteur propre de $A$ ssi
    c'est un vecteur propre de l'endomorphisme $u_A$, c'est à dire s'il existe 
    $\lambda \in \K$ tel que $AX = \lambda X$ et $X \neq 0$.
\end{definition}

\begin{definition}{Espace propre d'une matrice}{}
    \index{Espace propre!d'une matrice}
    Soit $\lambda \in \K$ une valeur propre de $A$. On appelle \emph{espace propre} 
    associé l'espace vectoriel 
    \begin{align} E_\lambda(A) & = E_\lambda(u_A) \\
                        & = \{ X \in \K^n \mid AX = \lambda X \} \\
                        & = \Ker (u_A - \lambda \Id_{\K^n}) 
    \end{align}
\end{definition}

\begin{definition}{Spectre d'une matrice}{}
    \index{Spectre!d'une matrice}
    L'ensemble des valeurs propres de $A$ est appelé spectre de $A$ 
    et noté $\Sp(A)$.
    
\end{definition}

\begin{definition}{Équation aux éléments propres d'une matrice}{}
    \index{Equation aux elements propres@\'Equation aux éléments propres!d'une matrice}
    On appelle \emph{équation aux éléments propres} pour $A$ l'équation 
    $$AX = \lambda X$$

    \begin{itemize}
        \item $\lambda \in K$
        \item $X \in \K^n \setminus \{0\}$
    \end{itemize}
    
\end{definition}

\begin{propositionnt}{lien entre éléments propres d'un endomorphisme et de sa matrice dans une base}{}
    Soit $E$ un $\K$-espace vectoriel de dimension finie $n$, 
    $B = (e_1, \ldots, e_n)$ une base de $E$ et $u \in \mathscr L(E)$, 
    $A = M_B(u) \in \mathscr M_n(\K)$. 
    Pour $x \in E$, on note $X = M_B(x)$ le vecteur colonne de ses coordonnées dans la base $B$, et alors 
    $\fonction{\varphi}{E}{\K^n}{x}{X}$ est un isomorphisme. 
    
    On a donc : 
    \begin{itemize}
        \item $\forall \lambda \in \K$, $u(x) = \lambda x$ ssi $AX = \lambda X$ donc $\lambda$ est une valeur propre de $u$
        ssi $\lambda$ est une valeur propre de $A$
        \item $x$ est un vecteur propre de $u$
        ssi $X$ est un vecteur propre de $A$.
    \end{itemize}
    Donc en particulier, $\Sp(u) = \Sp(A)$.
\end{propositionnt}


\begin{proposition}{Matrices semblables et spectre}{}
    Deux matrices semblables ont le même spectre.
\end{proposition}

\idee{
    Changement de base.
}

\begin{proof}
    Soient $A, A' \in \mathscr M_n(\K)$ deux matrices semblables, et 
    soit $u \in \mathscr L(\K^n)$ canoniquement associée à $A$. 

    Or $A' = P^{-1} A P$ pour une certaine matrice inversible $P \in \mathscr M_n(\K)$, 
    donc $A'$ est la matrice de $u$ dans une autre base de $\K^n$.

    Par ce qui précède, $A$ et $A'$ ont le même spectre.
\end{proof}

\begin{proposition}{Spectre selon un sous-corps}{}
    Soit $\K \subset \K'$ sous-corps de $\K'$. 
    Soit $A \in \mathscr M_n(\K)$. Alors on peut voir $A$ comme une matrice de 
    $\mathscr M_n(\K')$ et les endomorphismes 
    $\fonction{u_A}{\K^n}{\K^n}{X}{AX}$ et
    $\fonction{u'_A}{\K'^n}{\K'^n}{X}{AX}$ peut avoir 2 spectres différents 
    $\Sp_{\K} (A)$ et $\Sp_{\K'}(A)$.

    Alors $\Sp_{\K} (A) \subset \Sp_{\K'}(A) \cap \K$.
\end{proposition}

\begin{proof}
    Soit $\lambda \in \Sp_{\K} (A)$ alors il existe $X \in \K^n \setminus \{0\}$ tel que
    $AX = \lambda X$.

    Or $\K \subset \K'$ donc $\lambda \in \K'$ et $X \in \K'^n \setminus \{0\}$ donc
    $\lambda \in \Sp_{\K'}(A)$.
\end{proof}

\begin{exemple}
    $A = \begin{pmatrix}
        0 & -1 \\
        1 & 0
    \end{pmatrix}$ alors $\Sp_{\R}(A) = \varnothing$ 

    Mais $\Sp_{\C}(A) = \{i, -i\}$.
\end{exemple}


\section{Polynôme caractéristique }
\subsection{Polynôme caractéristique d'une matrice carrée}

\begin{definition}{Polynôme caractéristique d'une matrice}{}
    \index{Polynôme caractéristique!d'une matrice}
    Soit $A \in \mathscr M_n (\K)$. On définit son polynôme caractéristique par : 
    \[
        \chi_A(X) = \sum_{\sigma \in S_n} \varepsilon(\sigma) \prod_{i = 1 }^{ n } (X \delta_{\sigma(i)}^i - a_{\sigma(i), i})
    \]

    Alors $\forall x \in \K$, $\chi_A(x) = \det (x I_n - A)$.

    On utilisera donc l'abus de notation $\chi_A(X) = \det (X I_n - A)$.
\end{definition}

\begin{theoreme}{Propriétés du polynôme caractéristique}{}
    Soit $A \in \mathscr M_n (\K)$, alors $\deg \chi_A = n$, $\chi_A$ est unitaire et en notant $\displaystyle \chi_A = \sum_{k = 0 }^{n } a_kX^k$ on a : 

    \begin{itemize}
        \item $a_n = 1$
        \item $a_{n - 1} = - \tr(A)$
        \item $a_0 = (-1)^n \det(A)$
    \end{itemize}
\end{theoreme}

\begin{exemple}
    $A = \begin{pmatrix}
        1 & 2 \\
        3 & 4
    \end{pmatrix}$ et donc $\chi_A = \begin{vmatrix}
        X - 1 & -2 \\
        -3 & X - 4
    \end{vmatrix} (X - 1)(X - 4) - 6 = X^2 - 5X - 2$.
\end{exemple}

\idee{
    Séparer le polynôme en deux parties : la partie diagonale avec $\sigma = \Id$ et le reste.
}

\begin{proof}
    Par le formule, $\chi_A$ est un polynôme de degré au plus $n$.

    $a_0 = \chi_A(0) = \det(-A) = (-1)^n \det(A)$.

    \avspace Posons : 
    
    $\displaystyle P_A = \prod_{i = 1 }^{ n } (X - a_{i, i})$ est un polynôme unitaire de degré $n$ scindé.

    $\displaystyle Q_A = \sum_{\sigma \in S_n \setminus \{\Id\}} \varepsilon(\sigma) \prod_{i = 1 }^{ n } (X \delta_{\sigma(i)}^i - a_{\sigma(i), i})$ 

    Soit $\sigma \in S_n$ tel que $\sigma \neq \Id$, et 
    $\displaystyle Z_{A, \sigma} = \varepsilon(\sigma) \prod_{i = 1 }^{ n } (X \delta_{\sigma(i)}^i - a_{\sigma(i), i})$.

    Or $\sigma \neq \Id$ donc $\exists i_0 \in \ninter 1 n$ tel que $\sigma(i_0) \neq i_0$ donc $\delta_{\sigma(i_0)}^{i_0} = 0$.

    Posons $i_1 = \sigma(i_0)$. On a déjà $i_1 \neq i_0$ et $\sigma$ injective, donc $\sigma(i_1) \neq \sigma(i_0)$
    donc $\sigma(i_1) \neq i_1$. 

    Donc dans le produit, il y a au moins deux indices $i_0$ et $i_1$ tels que 
    $\delta_{\sigma(i)}^i = 0$ donc le degré de $Z_{A, \sigma}$ est au plus $n - 2$.

    Donc $\deg Q_A \leqslant n - 2$, donc les coefficients de $X^n$ et $X^{n - 1}$ dans $\chi_A$ 
    sont ceux de $P_A$ donc valent $1$ et $\displaystyle - \sum_{i = 1 }^{n } a_{i, i} = - \tr(A)$.
\end{proof}

\begin{proposition}{}{}
    Soit $A \in \mathscr M_n (\K)$ matrice triangulaire, $A = (a_{i, j})$. Alors 
    $\displaystyle \chi_A(X) = \prod_{i = 1 }^{ n } (X - a_{i, i})$.
\end{proposition}

\begin{proof}
    Si $A$ est triangulaire supérieure, alors 
    $$\chi_A = \begin{vmatrix}
        X - a_{1, 1} & -a_{1, 2} & \cdots & -a_{1, n} \\
        0 & X - a_{2, 2} & \cdots & -a_{2, n} \\
        \vdots & \vdots & \ddots & \vdots \\
        0 & 0 & \cdots & X - a_{n, n}
    \end{vmatrix} = \prod_{i = 1 }^{ n } (X - a_{i, i})$$
\end{proof}

\begin{definition}{Matrice compagnon, HP}{}
    Soit $P$ polynôme unitaire de degré $n$. Notons 
    $\displaystyle P = X^n - \sum_{k = 0 }^{n - 1} a_k X^k$.

    Notons $A = \begin{pmatrix}
        0 & 0 & \cdots & 0 & a_0 \\
        1 & 0 & \cdots & 0 & a_1 \\
        0 & 1 & \cdots & 0 & a_2 \\
        \vdots & \vdots & \ddots & \vdots & \vdots \\
        0 & 0 & \cdots & 1 & a_{n - 1}
    \end{pmatrix} \in \mathscr M_n(\K)$ la matrice compagnon de $P$. On a donc une matrice par blocs, 
    et on a alors $\chi_A = P$

\end{definition}

\idee{
    Développement par rapport à la dernière colonne.
}

\begin{proof}
    $\chi_A = \det(X I_n - A ) = \begin{vmatrix}
        X & 0 & 0 & \dots & -a_0\\
        -1 & 0 & 0 & & -a_1 \\
        0 & \ddots & & \ddots &  \\
        \vdots & & \ddots & & \vdots \\
        0 & \dots & 0 & -1 & X - a_{n - 1}
    \end{vmatrix}$

    Développons par rapport à la dernière colonne : 

    $\chi_A = (-1)^{n + 1} (- a_0) \begin{vmatrix}
        -1 & \dots & \star \\
        \vdots & \ddots & \vdots \\
        0 & \dots & -1
    \end{vmatrix} + \displaystyle \sum_{i = 1 }^{ n - 1 } (-1)^{n + i} (-a_{i - 1}) 
    \begin{vmatrix}
        X & & 0 &  & & \\
         & \ddots & &  & 0 &  \\
        * & & X &  & & \\
         &  & & -1 & & *\\
         & 0 & & & \ddots & \\
         &  &  & 0 & & -1
    \end{vmatrix} + (-1)^{2n} (X - a_{n - 1}) 
    \begin{vmatrix}
        X &  &  & 0 \\
         & X &  &  \\
         & & \ddots & \\
        * & &  & X
    \end{vmatrix}$

    $\displaystyle \chi_A = (-1)^n a_0 (-1)^{n - 1} + \sum_{i = 2 }^{ n -1} (-1)^{n + i + 1} X^{i - 1} (-1)^{n - i} a_{i - 1} + (X - a_{n - 1}) X^{n - 1}$

    $\displaystyle \chi_A = - a_0 + \sum_{i = 2 }^{ n -1} (-a_{i - 1}) X^{i - 1} - a_{n - 1} X^{n - 1} + X^n = P$
    
\end{proof}

\begin{exemple}
    $A = \begin{pmatrix}
        0 & 0 & 1 \\
        1 & 0 & 2 \\
        0 & 1 & 3
    \end{pmatrix}$ et donc $\chi_A = \begin{vmatrix}
        X & 0 & -1 \\
        -1 & X & -2 \\
        0 & -1 & X - 3
    \end{vmatrix} = X^3 - 3X^2 + 2X - 1$.
\end{exemple}

\subsection{Polynôme caractéristique d'un endomorphisme}

\begin{theoreme}{Polynôme caractéristique d'un endomorphisme}{}
    \index{Polynôme caractéristique!d'un endomorphisme}
    Soit $E$ un $\K$-espace vectoriel de dimension finie $n$ et $u \in \mathscr L(E)$.

    Soit $B$ une base de $E$ et $A = M_B(u) \in \mathscr M_n(\K)$. 

    Alors $\chi_A$ ne dépend pas du choix de la base $B$.

    On l'appelle \emph{polynôme caractéristique} de $u$ et on le note $\chi_u = \chi_A$.

    On notera de même $\chi_u(X) = \det(X \Id_E - u)$.
\end{theoreme}

\idee{
    Astuce du topin.
}

\begin{proof}
    Soit $B'$ une autre base, et $A' = M_{B'}(u)$.

    Alors $\exists P \in GL_n(\K)$ tel que $A = P A' P^{-1}$.

    Pour $x \in \K$, 
    \begin{align}
        \chi_A(x) & = \det(x I_n - A) \\
                        & = \det(x PP^{-1} - P A' P^{-1}) \\
                        & = \det(P(x I_n - A') P^{-1}) \\
                        & = \det(P) \det(x I_n - A') \det(P^{-1}) \\
                        & = \chi_{A'}(x)
    \end{align}

    Donc $\chi_A = \chi_{A'}$.
\end{proof}

\begin{corollaire}{matrices semblables et polynôme caractéristique}{}
    Deux matrices semblables ont le même polynôme caractéristique.
\end{corollaire}

\begin{proof}
    Résulte directement de la preuve du théorème précédent.
\end{proof}


\subsection{Lien avec les valeurs propres }

\begin{theoreme}{Valeurs propres et polynôme caractéristique}{}
    Les valeurs propres d'un endomorphisme en dimension finie (respectivement d'une matrice carrée) 
    sont les racines de son polynôme caractéristique.
\end{theoreme}

\begin{proof}
    Soit $u \in \mathscr L(E)$ et $\lambda$ valeur propre de $u$ 
    ssi $\ker(u - \lambda \Id_E) \neq \{0\}$
    ssi $u - \lambda \Id_E \notin GL(E)$
    ssi $\det(u - \lambda \Id_E) = 0$
    ssi $\chi_u(\lambda) = 0$.
\end{proof}

\begin{exemple}
    $E = \{f \in \mathscr C^\infty(\R, \R), f'' - 3 f' + 2f = 0\}$ $E$ est un 
    $\R$-espace vectoriel de dimension $2$.

    Considérons $u : f \mapsto f'$ endomorphisme de $E$ car linéaire. 

    Si $f \in E$, $f'$ est $\mathscr C^\infty$ et $f'' - 3 f' + 2f = 0$ donc $f''' - 3 f'' + 2 f' = 0$ donc $f' \in E$.

    Remarquons que $f_1(x) = e^{2x} - e^x$ et $f_2(x) = e^{2x} + e^x$ sont dans $E$. 

    $u(f_1) = f_1 + f_2 - \frac{f_2 - f_1}{2} = \frac{3}{2} f_1 + \frac{1}{2} f_2$ et
    $u(f_2) = f_2 + f_1 + \frac{f_2 - f_1}{2} = \frac{1}{2} f_1 + \frac{3}{2} f_2$.

    $M_{f_1, f_2}(u) = \begin{pmatrix}
        \frac{3}{2} & \frac{1}{2} \\
        \frac{1}{2} & \frac{3}{2}
    \end{pmatrix}$.

    Donc $\chi_u = \begin{vmatrix}
        X - \frac{3}{2} & -\frac{1}{2} \\
        -\frac{1}{2} & X - \frac{3}{2}
    \end{vmatrix} = X^2 - 3X + 2 = (X - 1)(X - 2)$.

    Donc $\Sp(u) = \{1, 2\}$.

    $g_1 : x \mapsto e^x$ est vecteur propre associé à la valeur propre $1$ car $g_1' = g_1$.

    $g_2 : x \mapsto e^{2x}$ est vecteur propre associé à la valeur propre $2$ car $g_2' = 2 g_2$.

    Donc $g_1, g_2$ est libre donc base de $E$.
\end{exemple}

\begin{exemple}
    $E = \{(u_n)_{n \in \N} \in \C^{\N}\}$ $\forall n \in \N u_{n + 2} = 4 u_{n + 1} - 3 u_n$.

    $E$ est un $\C$-espace vectoriel. 

    $\fonction \varphi E {\C^2} {(u_n)_{n \in \N}} {(u_0, u_1)}$ est un isomorphisme.

    $b_1 = (1, 0)$ et $b_2 = (0, 1)$ est une base de $\C^2$.

    Donc $v = \varphi^{-1} (b_1) $ et $w = \varphi^{-1} (b_2)$ est une base de $E$.

    Considérons $s : (u_n)_{n \in \N} \mapsto (u_{n + 1})_{n \in \N}$ endomorphisme de $E$.

    Cherchons sa matrice dans la base $(v, w)$ avec $\varphi(v) = (1, 0)$ donc 
    $v_0 = 1$ et $v_1 = 0$    
    et $\varphi(w) = (0, 1)$ donc $w_0 = 0$ et $w_1 = 1$.

    Posons $v' = s(v) = (v_{n + 1})_{n \in \N}$ et $w' = s(w) = (w_{n + 1})_{n \in \N}$.

    On a alors $v' = -3w$ et $w' = v + 4w$.

    Donc $M_{(v, w)}(s) = \begin{pmatrix}
        0 & 1 \\
        -3 & 4
    \end{pmatrix}$.

    Donc $\chi_s = \begin{vmatrix}
        X & -1 \\
        3 & X - 4
    \end{vmatrix} = X^2 - 4X + 3 = (X - 1)(X - 3)$.

    Donc $\Sp(s) = \{1, 3\}$.


    $u = (u_n)_{n \in \N}$ est vecteur propre associé à la valeur propre $1$ ssi 
    $s(u) = u$ ssi $\forall n \in \N, u_{n + 1} = u_n$ ssi $\forall n, u_n = \alpha$. 

    $u$ est vecteur propre pour $3$ ssi $\forall n, u_{n + 1} = 3 u_n$ ssi $\forall n, u_n = \beta 3^n$

    $((1), (3^n))$ sont deux vecteurs propres pour des valeurs propres différentes donc famille libre donc base de 
    $E$. 
\end{exemple}


\begin{proposition}{}{}
    Soit $u \in \mathscr L(E)$ et $F$ sous-espace vectoriel de $E$ stable par $u$, et $\tilde u$
    l'endomorphisme induit par $u$ sur $F$. 

    Alors 
    $\chi_{\tilde u} \mid \chi_u$ et en particulier, $\Sp(\tilde u) \subset \Sp (u)$
\end{proposition}

\idee{Par blocs}

\begin{proof}
    Soit $B$ une base de $F$ de dimension $p$ et $B = B_1 \sqcup B_2$ base de $E$, alors 
    $M_B(u) = \begin{pmatrix}
        M_{B_1}(\tilde u) & C \\
        0 & D
    \end{pmatrix}$.

    \begin{align}
        \chi_u(X) & = \det(X I_n - M_B(u)) \\
                    & = \begin{vmatrix}
                        X I_p - M_{B_1}(\tilde u) & -C \\
                        0 & X I_{n - p} - D
                    \end{vmatrix} \\
                    & = \det(X I_p - M_{B_1}(\tilde u)) \det(X I_{n - p} - D) \\
                    & = \chi_{\tilde u}(X) Q (X)
    \end{align}
\end{proof}

\begin{definition}{Multiplicité d'une valeur propre}{}
    \index{Multiplicité!d'une valeur propre}
    Soit $u \in \mathscr L(E)$ et $\lambda \in \Sp(u)$,
    alors $\lambda $ est racine de $\chi_u$.

    On appelle ordre de multiplicité de la valeur propre $\lambda$ son ordre de multiplicité en tant que racine de $\chi_u$.
\end{definition}

\begin{exemple}
    Une valeur propre de multiplicité 1 sera dite valeur propre simple, double pour multiplicité 2, etc.

    On note souvent l'ordre $m_\lambda$. On a la même notion pour les matrices carrées. 
\end{exemple}

\begin{proposition}{Dimension sous-espace propre}{}
    Soit $u \in \mathscr L(E), \, E$ de dimension finie et $\lambda \in \Sp(u)$. 

    Alors en notant $m_\lambda$ son ordre de multiplicité et $E_\lambda(u)$ son espace propre associé, on a :
    $$1 \leqslant \dim(E_\lambda(u)) \leqslant m_\lambda$$

    En particulier si $\lambda$ est une valeur propre simple, $\dim E_\lambda (u) = 1$.
\end{proposition}

\begin{proof}
    $E_\lambda(u)$ est stable par $u$. L'endomorphisme induit est $\tilde u = \lambda \Id_{E_\lambda(u)}$
    or $\chi_{\tilde u} \mid \chi_u$ et 
    $\chi_{\tilde u} = \begin{vmatrix}
        X - \lambda & 0 & \cdots & 0 \\
        0 & X - \lambda & \cdots & 0 \\
        \vdots & \vdots & \ddots & \vdots \\
        0 & 0 & \cdots & X - \lambda
    \end{vmatrix} = (X - \lambda)^p$

    Donc $(X - \lambda)^p \mid \chi_u$ donc $p \leqslant m_\lambda$.
\end{proof}

\begin{remarque}
    Pour calculer $\dim E_\lambda(u)$ on peut utiliser le théorème du rang. 

    $\dim E_\lambda(u) = \dim \Ker (u - \lambda \Id_E) = \dim E - \rg(u - \lambda \Id_E)$.
\end{remarque}



\begin{exemple}
    $A = \begin{pmatrix}
        8 & 2 & 2 \\
        2 & 8 & 2 \\
        2 & 2 & 8
    \end{pmatrix}$.

    Alors $A \cdot \begin{pmatrix}
        1 \\
        1 \\
        1
    \end{pmatrix} = 12 \cdot \begin{pmatrix}
        1 \\
        1 \\
        1
    \end{pmatrix}$ donc $12$ est valeur propre, et $\begin{pmatrix}
        1 \\
        1 \\
        1
    \end{pmatrix}$ vecteur propre de $A$. Donc $\dim E_{12}(A) \geqslant 1$.

    De plus $A - 6 I_3 = \begin{pmatrix}
        2 & 2 & 2 \\
        2 & 2 & 2 \\
        2 & 2 & 2
    \end{pmatrix}$ donc $\rg(A - 6 I_3) = 1$ donc $\dim E_6(A) = \dim \ker (A - 6 I_3) = 3 - 1 = 2$.

    Donc $6$ est valeur propre de multiplicité au moins $2$. Or $E_{12}(A) + E_6(A)$ est directe, donc 
    $\dim E_{12}(A) + \dim E_6(A) \leqslant 3$ donc $E = E_{12}(A) \oplus E_6(A)$ et $\dim E_{12}(A) = 1$.

    De plus $m_{12} \geqslant 1$ et $m_6 \geqslant 2$ et $\deg \chi_A = 3$ donc $m_{12} = 1$ et $m_6 = 2$
    et $\chi_A = (X - 12)(X - 6)^2$.

\end{exemple}


\begin{proposition}{Polynôme caractéristique scindé}{}
    Comme $\K \subset \C$, on pourra toujours se ramener au cas de $\chi_A$ complexe, en 
    considérant les valeurs propres complexes.

    Soit $A \in \mathscr M_n(\K)$, alors $\Sp(A) \neq \emptyset$ car $\chi_A$ admet au moins une racine dans $\K$
    et en notant $\Sp(A) = \{\lambda_1, \ldots, \lambda_p \}$, $\chi_A$ est scindé et on a :
    $$\chi_A(X) = \prod_{i = 1 }^{ p } (X - \lambda_i)^{m_{\lambda_i}}$$

    Donc par les relations coefficients racines, 
    \begin{itemize}
        \item $\displaystyle \tr(A) = \sum_{i = 1 }^{ p } m_{\lambda_i} \lambda_i$ 
        \item $\displaystyle \det A = \prod_{i = 1 }^{ p } \lambda_i^{m_{i}}$
    \end{itemize}
\end{proposition}

\begin{exemple}
    $A = \begin{pmatrix}
        1 & 7 & 3 \\
        -1 & 4 & 8 \\
        3 & 10 & -2
    \end{pmatrix}$. On cherche le polynôme caractéristique, valeurs propres et vecteurs propres. 

    \uindent{Solution 1 } 

    \svspace On calcul $\chi_A = \begin{vmatrix}
        X - 1 & -7 & -3 \\
        1 & X - 4 & -8 \\
        -3 & -10 & X + 2
    \end{vmatrix}$ et on cherche ses racines. Puis pour chaque racine $\lambda$, on cherche le rang de 
    $(A - \lambda I_3)$.

    \uindent{Solution 2 }

    \svspace On remarque que $L_2 + L_3 = 2 L_1$ donc $A$ non inversible, donc $0$ est valeur propre de $A$. 

    $A \cdot \begin{pmatrix}
        1 \\
        1 \\
        1
    \end{pmatrix} = 11 \cdot \begin{pmatrix}
        1 \\
        1 \\
        1
    \end{pmatrix}$ donc $11$ est valeur propre de $A$.

    Donc $\chi_A$ est divisble par $X(X - 11)$.

    Donc il vient $\chi_A(X) = X(X - 11)(X - a)$ avec $a \in \R$.

    $\tr A = 3 = 11 + a$ donc $a = -8$.

    Donc $\chi_A(X) = X(X - 11)(X + 8)$, donc $\Sp(A) = \{0, 11, -8\}$ sont valeurs propres simples.

    Les espaces propres sont donc de dimension 1. Cherchons leurs bases. 

    Soit $U = \begin{pmatrix}
        x \\
        y \\
        z   
    \end{pmatrix}$ alors $U \in E_{-8}(A)$ ssi $(A + 8I) U = 0$ ssi $\begin{cases}
        (1) 9x + 7y + 3z = 0 \\
        (2) -x + 12y + 8z = 0 \\
        (3) 3x + 10y + 6z = 0
    \end{cases}$ ssi $\begin{cases}
        -x +  & 12y + 8z = 0 (2)\\
        & 115y + 75 z = 0 (1) + 9*(2) \\
        & 46 y + 30 z = 0 (3) - 3*(2)
    \end{cases}$ 
    
    Et les dernières équations sont proportionnelles. 

    Posons $e_{-8} = \begin{pmatrix}
        4 \\ 
        -15 \\
        23
    \end{pmatrix}$ Or $\dim E_{-8}(A) = 1$ donc $E_{-8}(A) = \Vect(e_{-8})$.

    De même pour les deux autres. 
\end{exemple}

\section{Endomorphismes et matrices carrées diagonalisables}

\begin{definition}{Endomorphisme diagonalisable}{}
    \index{Endomorphisme!diagonalisable}
    Soit $u \in \mathscr L(E)$ où $E$ est de dimension finie. On dit que $u$ est diagonalisable ssi il existe 
    $B$ une base de $E$ telle que $M_B(u)$ est une matrice diagonale ssi il existe 
    $B$ base de $E$ formée de vecteurs propres de $u$.

    En ce cas, $M_B(u) = \begin{pmatrix}
        \lambda_1 & 0 & \cdots & 0 \\
        0 & \lambda_2 & \cdots & 0 \\
        \vdots & \vdots & \ddots & \vdots \\
        0 & 0 & \cdots & \lambda_n
    \end{pmatrix}$ et les éléments diagonaux sont les valeurs propres de $u$, répétées avec multiplicité car 
    $\displaystyle \chi_u = \prod_{i = 1 }^{ n } (X - \lambda_i)$.
\end{definition}

\begin{exemple}
    Prenons $u$ une projection. Supposons que $E = F \oplus G$ avec $F = \im(u)$ et $G = \Ker(u)$.

    Soit $B = (e_1, \dots, e_p, e_{p + 1}, \dots, e_n)$ une base de $E$ telle que 
    $(e_1, \dots, e_p)$ base de $F$ et $(e_{p + 1}, \dots, e_n)$ base de $G$.

    $M_B(u) = \begin{pmatrix}
        I_p & 0 \\
        0 & 0
    \end{pmatrix}$ donc $u$ est diagonalisable.

    Donc $\chi_u(X) = (X - 1)^p X^{n - p}$. $1$ est valeur propre de multiplicité $p$, espace propre 
    $F$ et $0$ de multiplicité $n - p$, espace propre $G$.

    \lvspace Pour $u$ une symétrie, $M_B(u) = \begin{pmatrix}
        I_p & 0 \\
        0 & -I_{n - p}
    \end{pmatrix}$ donc $u$ est diagonalisable.
\end{exemple}

\subsection{Caractérisation}

\begin{theoreme}{Caractérisation d'un endomorphisme diagonalisable}{}
    Soit $u \in \mathscr L(E)$ et $\Sp(u) = \{\lambda_1, \ldots, \lambda_p\}$.

    Alors $u$ est diagonalisable ssi $E$ est la somme des sous-espaces propres de $u$ :
    $$E = \bigoplus_{j = 1 }^{ p } E_{\lambda_j}(u)$$
\end{theoreme}

\idee{
    Choisir la bonne base.
}

\begin{proof}
    $\Rightarrow$ Soit $B = (e_1, \ldots, e_n)$ base de $E$ telle que $M_B(u)$ est diagonale.

    Pour $i \in \ninter 1 n$, $e_i$ est un vecteur propre de $u$, donc $\exists j \in \ninter 1 p$
    tel que $e_i \in E_{\lambda_j}(u)$, donc $e_i \in \bigoplus_{j = 1 }^{ p } E_{\lambda_j}(u)$.

    Donc $E = \Vect(e_1, \ldots, e_n) \subset \bigoplus_{j = 1 }^{ p } E_{\lambda_j}(u)$.

    \lvspace $\Leftarrow$ Supposons que $\displaystyle E = \bigoplus_{j = 1 }^{ p } E_{\lambda_j}(u)$.

    Soit $B$ base de $E$ adaptée à cette somme directe. Chaque vecteur de $B$ est propre donc 
    $u$ est diagonalisable.
\end{proof}

\begin{corollairent}{caractérisation dimensionnelle d'un endomorphisme diagonalisable}{}
    Soit $u \in \mathscr L(E)$ de spectre $\Sp(u) = \{\lambda_1, \ldots, \lambda_p\}$.
    Alors $u$ est diagonalisable ssi $\displaystyle \sum_{j = 1 }^{p } \dim E_{\lambda_j}(u) = \dim E$.
\end{corollairent}

\begin{proof}
    $u$ diagonalisable ssi $\displaystyle E = \bigoplus_{j = 1 }^{ p } E_{\lambda_j}(u)$
    
    ssi $\displaystyle n = \dim E = \sum_{j = 1 }^{ p } \dim E_{\lambda_j}(u)$.

    ssi $\displaystyle \sum_{j = 1 }^{ p } \dim E_{\lambda_j}(u) = n$.
\end{proof}


\begin{theoreme}{Caractérisation par les multiplicités}{}
    Soit $u \in \mathscr L(E)$ tel que $\Sp(u) = \{\lambda_1, \ldots, \lambda_p\}$.

    $u$ est diagonalisable ssi $\chi_u$ est scindé et $\forall \lambda \in \Sp(u), \, m_\lambda = \dim E_\lambda(u)$.
\end{theoreme}

\begin{proof}
    $\Rightarrow$ Supposons que $u$ est diagonalisable.

    Alors il existe $B$ base telle que $M_B(u) = \operatorname{Diag}(\lambda_1, \ldots, \lambda_p)$
    et $\displaystyle \chi_u(X) = \prod_{i = 1 }^{ p } (X - \lambda_i)^{m_j}$ est scindé. 

    Par ailleurs pour tout $j$, les $m_j$ vecteurs de la base associés à la valeur propre $\lambda_j$
    forment une base de $E_{\lambda_j}(u)$ donc $\dim E_{\lambda_j}(u) = m_j$. 

    \lvspace $\Leftarrow$ Supposons que $\chi_u$ est scindé et que $\forall j, \dim E_\lambda(u) = m_j$.

    Alors $\displaystyle \sum_{j = 1 }^{ p } m_j = n$ donc $\displaystyle \sum_{j = 1 }^{ p } \dim E_{\lambda_j}(u) = n$, 

    donc $u$ est diagonalisable.
\end{proof}


\begin{proposition}{Diagonalisabilité et polynôme caractéristique}{}
    Soit $u \in \mathscr L(E)$. Si $\chi_u$ est scindé à racines simples, alors $u$ est diagonalisable.
\end{proposition}

\begin{proof}
    En ce cas $\chi_u$ scindé et pour toute valeur propre $\lambda$, $1 \leqslant \dim E_\lambda(u) \leqslant m_\lambda = 1$
\end{proof}

\begin{remarque}
    Attention, la réciproque est fausse. Par exemple, $u = 0$ alors $M_B(u) = 0$ diagonalisable, pourtant 
    $\chi_u = X^n$ n'est pas scindé à racines simples.
\end{remarque}

\subsection{Diagonalisation d'une matrice carrée }

\begin{definition}{Matrice diagonalisable}{}
    \index{Matrice!diagonalisable}
    Soit $A \in \mathscr M_n(\K)$. On dit que $A$ est diagonalisable ssi $A$ est semblable à une matrice diagonale
    ssi $\exists P \in GL_n(\K), \, \exists D \in \mathscr M_n(\K)$ diagonale, telle que $A = P D P^{-1}$.
\end{definition}

\begin{proposition}{Lien diagonalisabilité matrice et endomorphisme}{}
    Soit $u \in \mathscr L(E)$, $B$ base de $E$ et $A = M_B(u)$. Alors $u$ est diagonalisable ssi $A$ est diagonalisable.
\end{proposition}

\begin{proof}
    $u$ est diagonalisable ssi il existe $B'$ base telle que $M_{B'}(u) = D$ est diagonale
    ssi $\exists P \in GL_n(\K)$ telle que $A = P D P^{-1}$ ($P$ matrice de passage de $B$ à $B'$).
\end{proof}

\begin{theoreme}{Théorème spectral}{}
    Toute matrice symétrique réelle $A$ est diagonalisable avec matrice de passage $P$ orthogonale
    c'est à dire telle que $P^{-1} = P^\top$.    
\end{theoreme}

\begin{proof}
    Plus tard dans l'année. 
\end{proof}

\begin{exemple}
    $A = \begin{pmatrix}
        3 & 2 & 2 \\
        2 & 3 & 2 \\
        2 & 2 & 3
    \end{pmatrix}$.

    1. \textbf{Valeur propre 7 :}
    On remarque que $V_1 = \begin{pmatrix} 1 \\ 1 \\ 1 \end{pmatrix}$ vérifie $AV_1 = 7V_1$.
    Donc $7$ est valeur propre et $\dim E_7(A) \geqslant 1$.

    2. \textbf{Valeur propre 1 :}
    Calculons $A - I = \begin{pmatrix}
        2 & 2 & 2 \\
        2 & 2 & 2 \\
        2 & 2 & 2
    \end{pmatrix}$.
    Cette matrice est de rang $1$ (toutes les colonnes sont colinéaires).
    D'après le théorème du rang, $\dim \ker(A - I) = 3 - \text{rang}(A-I) = 2$.
    Donc $\dim E_1(A) = 2$.

    3. \textbf{Conclusion :}
    La somme des dimensions des espaces propres est $1 + 2 = 3$, ce qui correspond à la taille de la matrice.
    On en déduit que $\dim E_7(A)$ est exactement $1$.
    Ainsi, la famille $(V_1)$ est bien une base de $E_7(A)$.

    \avspace Le polynôme caractéristique est donc $\chi_A(X) = (X - 7)(X - 1)^2$.

    $\begin{pmatrix}
        x \\
        y \\
        z
    \end{pmatrix} \in E_1(A)$ ssi $2x + 2y + 2z = 0$ ssi $x + y + z = 0$.

    Une base est $\begin{pmatrix}
        1 \\
        -1 \\
        0
    \end{pmatrix}, \begin{pmatrix}
        0 \\
        1 \\
        -1
    \end{pmatrix}$. Posons $D = \begin{pmatrix}
        7 & 0 & 0 \\
        0 & 1 & 0 \\
        0 & 0 & 1
    \end{pmatrix}$ et $P = \begin{pmatrix}
        1 & 1 & 0 \\
        1 & -1 & 1 \\
        1 & 0 & -1
    \end{pmatrix}$.

    Alors $A = P D P^{-1}$ et $P$ est inversible.
\end{exemple}



\subsection{Applications }

\begin{proposition}{Puissance d'une matrice diagonalisable}{}
    Soit $A \in \mathscr M_n(\K)$ diagonalisable. Notons $A = PDP^{-1}$ avec $P \in GL_n(\K)$ et $D$ diagonale.
    Alors par récurrence, $$\forall k \in \N, \, A^k = P D^k P^{-1}$$
\end{proposition}

\begin{proposition}{Suites récurrentes et matrices}{}
    Soient $(u_n)_{n \in \N}, (v_n)_{n \in \N}$ et $(w_n)_{n \in \N}$ définies par leur 
    premier terme et une relation de récurrence du type 
    $$\forall n \in \N : \begin{cases}
        u_{n + 1} = \alpha u_n + \beta v_n + \gamma w_n \\
        v_{n + 1} = \alpha' u_n + \beta' v_n + \gamma' w_n \\
        w_{n + 1} = \alpha'' u_n + \beta'' v_n + \gamma'' w_n
    \end{cases}$$

    Alors en posant $A = \begin{pmatrix}
        \alpha & \beta & \gamma \\
        \alpha' & \beta' & \gamma' \\
        \alpha'' & \beta'' & \gamma''
    \end{pmatrix}$ et $X_n = \begin{pmatrix}
        u_n \\
        v_n \\
        w_n
    \end{pmatrix}$, on a $X_{n + 1} = A X_n$ et donc $X_n = A^n X_0$.
\end{proposition}


\begin{exemple}
    \begin{tikzpicture}[scale=1.5, >=stealth]

        % --- Fond (Facultatif, pour simuler le tableau si désiré, ici blanc pour clarté) ---
        % \fill[black!80] (-4,-3) rectangle (4,4); 
        
        % --- Axes (Dessinés finement en gris comme sur le tableau) ---
        \draw[gray!50, thin] (-3.5,0) -- (3.5,0); % Axe X
        \draw[gray!50, thin] (0,-2.5) -- (0,3.5); % Axe Y

        % --- Définition des points (Coordonnées approximatives basées sur l'image) ---
        \coordinate (A) at (-2.5, 0.8);  % M0 = A
        \coordinate (B) at (1.8, 2.8);   % B = M1
        \coordinate (C) at (2.2, -1.8);  % C = M2

        % --- Calcul des points suivants (Barycentres) ---
        % M3 est le barycentre de (M0, M1, M2) aka (A, B, C)
        \coordinate (M3) at (barycentric cs:A=1,B=1,C=1);
        
        % M4 est le barycentre de (M1, M2, M3)
        \coordinate (M4) at (barycentric cs:B=1,C=1,M3=1);

        % --- Tracé du Triangle Principal (M0-M1-M2) ---
        \draw[thick] (A) -- (B) -- (C) -- cycle;

        % --- Tracé des lignes intérieures (Suite récursive) ---
        % Lignes vers M3
        \draw[thin] (A) -- (M3);
        \draw[thin] (B) -- (M3);
        \draw[thin] (C) -- (M3);

        % --- Tracé de l'étape suivante (M4) en Rouge ---
        % Sur la photo, on voit un triangle formé à droite et une croix rouge
        \draw[red, thin] (B) -- (M4) -- (C);
        \draw[red, thin] (M3) -- (M4);
        
        % Croix rouge pour M4
        \node[red] at (M4) {$\times$};

        % --- Étiquettes (Labels) ---
        \node[left] at (A) {$M_0=A$};
        \node[above right] at (B) {$B=M_1$};
        \node[right] at (C) {$C=M_2$};
        
        % M3 (placé légèrement à gauche du point pour lisibilité)
        \node[left, font=\small, yshift=2pt] at (M3) {$M_3$};
        
        % M4 (en rouge et petit)
        \node[right, red, font=\small, xshift=2pt] at (M4) {$M_4$};

    \end{tikzpicture}


    $\forall n, z_{n + 3} = \frac 1 3 (z_n + z_{n + 1} + z_{n + 2})$, quelle est la limite 
    de $z_n$ si elle existe ?

    Posons $\forall n, \, X_n = \begin{pmatrix}
        z_n \\
        z_{n + 1} \\
        z_{n + 2}
    \end{pmatrix}$ donc $X_{n + 1} = \begin{pmatrix}
        0 & 1 & 0 \\
        0 & 0 & 1 \\
        \frac{1}{3} & \frac{1}{3} & \frac{1}{3}
    \end{pmatrix} X_n$.


    $\forall n, X_n = A^n X_0$ par récurrence. 

    $\chi_A = \begin{vmatrix}
        X & -1 & 0 \\
        0 & X & -1 \\
        -\frac{1}{3} & -\frac{1}{3} & X - \frac{1}{3}
    \end{vmatrix} = (X - 1)(X^2 + \frac{2}{3} X + \frac{1}{3})$.

    Racines complexes conjuguées pour le second polynôme, 
    
    donc $\ds \Sp(A) = \{1, \frac{-1 + i \sqrt{2}}{3}, \frac{-1 - i \sqrt{2}}{3}\}$.

    $D = \begin{pmatrix}
        1 & 0 & 0 \\
        0 & \frac{-1 + i \sqrt{2}}{3} & 0 \\
        0 & 0 & \frac{-1 - i \sqrt{2}}{3}
    \end{pmatrix}$, et posons $\ds \alpha = \frac{-1 + i \sqrt{2}}{3}$ et $\ds \bar \alpha = \frac{-1 - i \sqrt{2}}{3}$.

    $A = PDP^{-1}$ et $U = \begin{pmatrix}
        x \\
        y \\
        z
    \end{pmatrix}$ valeur propre pour $\alpha$ ssi 
    $\begin{cases}
        y = \alpha x \\
        z = \alpha z 
        \frac 1 3 (x + y + z) = \alpha z
    \end{cases}$ 

    et donc $U = \begin{pmatrix}
        1 \\
        \alpha \\
        \alpha^2
    \end{pmatrix}$ convient. 

    $P = \begin{pmatrix}
        1 & 1 & 1 \\
        1 & \alpha & \bar \alpha \\
        1 & \alpha^2 & \bar \alpha^2
    \end{pmatrix}
    $ matrice de Vandermonde inversible.

    $X_n = P \begin{pmatrix}
        1 & 0 & 0 \\
        0 & \alpha^n & 0 \\
        0 & 0 & \bar \alpha^n
    \end{pmatrix} P^{-1} X_0$.

    $|\alpha| = |\bar \alpha| < 1$ donc $\alpha^n \to 0$.

    et après calcul on trouve $\ds L = \begin{pmatrix}
        \frac{z_0 + 2 Z_1 + 3 Z_2}{6} \\
        \frac{z_0 + 2 Z_1 + 3 Z_2}{6} \\
        \frac{z_0 + 2 Z_1 + 3 Z_2}{6}
    \end{pmatrix}$.
\end{exemple}

\subsubsection{Résolution de systèmes différentiels croisés }

On s'intéresse à $X \in \mathscr C^1(\R, \R^n)$ telle que $\forall t \in \R, \, X'(t) = A X(t)$ 
avec $A \in \mathscr M_n(\R)$, c'est-à-dire qu'en notant $X(t) = \begin{pmatrix}
x_1(t) \\
\vdots \\
x_n(t)
\end{pmatrix}$, on a $$\forall t \in \R, \,(S) \begin{cases}
\ds x_1'(t) & = \sum_{j = 1 }^{n } a_{1j} x_j(t) \\
& \vdots \\
\ds x_n'(t) & = \sum_{j = 1 }^{ n } a_{nj} x_j(t)
\end{cases}$$

et donc si on arrive à diagonaliser $A$, on peut écrire $A = PDP^{-1}$ avec $D$ diagonale
et alors en posant $Y(t) = P^{-1} X(t)$, on a $Y \in \mathscr C^1(\R, \R^n)$ et
$$\forall t \in \R, \, Y'(t) = P^{-1} X'(t) $$ car $U \mapsto P^{-1} U$ est linéaire et donc dérivable en chaque point.

Donc $(S) \iff Y'(t) = D Y(t)$ que l'on résout facilement. 

On en déduit $X(t) = P Y(t)$.

\subsubsection{Recherche du commutant d'une matrice }

Pour $A \in \mathscr M_n(\K)$, on s'intéresse à l'ensemble $\mathscr C(A) = \{M \in \mathscr M_n(\K), AM = MA\}$.

Si $A = PDP^{-1}$ avec $D$ diagonale, on a 
$\fonction{}{\mathscr C (A)}{\mathscr C(D)}{M}{P^{-1} M P}$ est un isomorphisme d'algèbres.

Il suffit donc de chercher $\mathscr C(D)$ en utilisant le fait que si 
$ND = DN$, les espaces propres de $D$ sont stables par $N$. Ainsi en rengeant dans l'ordre les valeurs propres
de $D$, $N \in \mathscr C(D)$ ssi $N$ est diagonale par blocs, avec des blocs de la dimension des espaces propres. 

\begin{exemple}
    $\begin{pmatrix}
        2 & 0 & 0 \\
        0 & 2 & 0 \\
        0 & 0 & 3
    \end{pmatrix}$ alors $ND = DN$ ssi $N = \begin{pmatrix}
        a & b & 0 \\
        c & d & 0 \\
        0 & 0 & e
    \end{pmatrix}$.
\end{exemple}

\subsubsection{Recherche des racines carrées d'une matrice }

Soit $A \in \mathscr M_n(\K)$. On cherche $B \in \mathscr M_n(\K)$ tel que $B^2 = A$.

Remarquons que si $B^2 = A$, alors $AB = B^2B = BB^2 = BA$, donc $B \in \mathscr C(A)$.

On cherche $B$ avec des coefficients dans $\mathscr C(A)$.

\section{Endomorphismes et matrices carrées trigonalisables}
\subsection{Définitions}

\begin{definition}{Endomorphisme trigonalisable}{}
    \index{Endomorphisme!trigonalisable}
    Un endomorphisme d'un $\K$-espace vectoriel est dit trigonalisble ssi il existe 
    $B$ une base de $E$ dans laquelle sa matrice est triangulaire supérieure.
\end{definition}

\begin{exemple}
    Tout endomorphisme diagonalisable est trigonalisable.

    Soit $u$ canoniquement associé à $A = \begin{pmatrix}
        1 & 1 \\
        0 & 1
    \end{pmatrix}$. Alors $u$ n'est pas diagonalisable car
    $\chi_u = (X - 1)^2$ et $\dim E_1(u) = 1 < 2$.
\end{exemple}

\begin{propositionnt}{drapeau complet stable par un endomorphisme trigonalisable}{}
    Soit $u \in \mathscr L(E)$ triagonalisable dans une base $B = (e_1, \ldots, e_n)$.

    Alors $e_1$ est vecteur propre de $u$, et plus généralement 
    $\forall k \in \ninter 1 n$, $F_k = \Vect(e_1, \ldots, e_k)$ est stable par $u$.

    On dit (HP) que $(F_1, \dots, F_n)$ forment un drapeau complet de sous-espaces vectoriels stables par $u$.
\end{propositionnt}

\begin{definition}{Matrice trigonalisable}{}
    \index{Matrice!trigonalisable}
    Une matrice carrée est trigonalisable ssi elle est semblable à une matrice triangulaire supérieure.
\end{definition}

\begin{proposition}{Lien trigonalisabilité matrice et endomorphisme}{}
    Soit $u \in \mathscr L(E)$, $B$ une base de $E$, et $A = M_B(u)$.

    Alors $u$ est trigonalisable ssi $A$ est trigonalisable.
\end{proposition}

\begin{proof}
    Identique à la preuve de la proposition sur la diagonalisabilité.
\end{proof}

\begin{theorement}{caractérisation trigonalisable par le polynôme caractéristique}{}
    Un endomorphisme est trigonalisable ssi son polynôme caractéristique est scindé. 

    Il en est de même pour les matrices carrées.
\end{theorement}

\idee{Par récurrence.}

\begin{proof}
    Avec le point de vu matriciel par récurrence sur la taille de la matrice.

    Soit $A \in \mathscr M_1(\K)$. Alors $A = (a)$ est triangulaire supérieure donc trigonalisable.

    HR : supposons le résultat vrai pour toute matrice de taille $\leqslant n$. 

    Soit $A \in \mathscr M_{n + 1}(\K)$ telle que $\chi_A$ soit scindé.

    Soit $\lambda$ une racine de $\chi_A$, et $e_1$ un vecteur propre associé à $\lambda$.

    Complétons $E_1$ en une base $B$. Alors dans $B$, l'endomorphisme associé à $A$ a pour matrice 
    $A' = \begin{pmatrix}
        \lambda & * & \cdots & * \\
        0 & & & \\
        \vdots & & A_1 & \\
        0 & & &
    \end{pmatrix}$ avec $A_1 \in \mathscr M_n(\K)$ qui est donc semblable à $A$. 

    $\chi_A = \chi_{A'} = (X - \lambda) \chi_{A_1}$ donc $\chi_{A_1}$ est scindé.

    Donc par hypothèse de récurrence, $A_1$ est trigonalisable, donc
    $A_1 = P T P^{-1}$ avec $P \in GL_n(\K)$ et $T$ triangulaire supérieure.

    Donc $A' = \begin{pmatrix}
        \lambda & * & \cdots & * \\
        0 & & & \\
        \vdots & & P T P^{-1} & \\
        0 & & &
    \end{pmatrix}$.

    Posons $Q = \begin{pmatrix}
        1 & 0 & \cdots & 0 \\
        0 & & & \\
        \vdots & & P & \\
        0 & & &
    \end{pmatrix} \in GL_{n + 1}(\K)$ car $P$ inversible, et 
    $Q^{-1} = \begin{pmatrix}
        1 & 0 & \cdots & 0 \\
        0 & & & \\
        \vdots & & P^{-1} & \\
        0 & & &
    \end{pmatrix}$.

    Donc $A$ est semblable à $Q^{-1} A' Q = \begin{pmatrix}
        \lambda & * & \cdots & * \\
        0 & & & \\
        \vdots & & T & \\
        0 & & &
    \end{pmatrix}$ qui est triangulaire supérieure.
\end{proof}

\begin{corollaire}{trigonalisabilité sur $\C$}{}
    Dans un $\C$-espace vectoriel de dimension finie, tout endomorphisme est trigonalisable.

    Toute matrice de $\mathscr M_n(\C)$ est trigonalisable.
\end{corollaire}


\begin{corollaire}{polynôme caractéristique d'une matrice trigonalisable}{}
    Si $A$ est trigonalisable, alors $A = P \begin{pmatrix}
        \lambda_1 & * & \cdots & * \\
        0 & \lambda_2 & \cdots & * \\
        \vdots & \vdots & \ddots & \vdots \\
        0 & 0 & \cdots & \lambda_n
    \end{pmatrix} P^{-1}$.

    Donc $\displaystyle \chi_A(X) = \prod_{i = 1 }^{ n } (X - \lambda_i)$ est scindé.
\end{corollaire}


\begin{exemple}
    $A = \begin{pmatrix}
        2 & -2 & -3 \\
        2 & -1 & -6 \\
        -1 & 2 & 4
    \end{pmatrix}$, et alors $A - I$ est de rang $1$. Donc $\dim \ker (A - I) = 2$, donc 
    $1$ est valeur propre de multiplicité $\leqslant 2$.

    Donc $\chi_A(X) = (X - 1)^2 (X - a)$ avec $a \in \R$.

    Or $\tr A = 3 = 2 + a$ donc $a = 1$ et donc $\chi_A(X) = (X - 1)^3$.

    Donc $1$ est valeur propre de multiplicité $3$, $A$ non diagonalisable
    mais $\chi_A$ est scindé donc $A$ est trigonalisable.

    \avspace Cherchons les vecteurs propres de $A$. 

    $\begin{pmatrix}
        x \\
        y \\
        z
    \end{pmatrix} \in E_1(A)$ ssi $x - 2y - 3z = 0$, donc $e_1 = \begin{pmatrix}
        2 \\
        1 \\
        0
    \end{pmatrix}, e_2 = \begin{pmatrix}
        3 \\
        0 \\
        1
    \end{pmatrix}$ forment une base de $E_1(A)$.

    Pour trigonaliser $A$, complétons en une base de $\R^3$ avec $e_3 = \begin{pmatrix}
        1 \\
        0 \\
        0
    \end{pmatrix}$.

    Alors $Ar_3 = \begin{pmatrix}
        2 \\
        2 \\
        -1
    \end{pmatrix} = 2 e_1 - e_2 + e_3$ 
    et donc dans la base $B = (e_1, e_2, e_3)$, l'endomorphisme associé à $A$ a pour matrice 
    $A' = \begin{pmatrix}
        1 & 0 & 2 \\
        0 & 1 & -1 \\
        0 & 0 & 1
    \end{pmatrix}$ triangulaire supérieure.
\end{exemple}


\subsection{Trace et déterminant }

\begin{proposition}{Trace et déterminant d'un endomorphisme trigonalisable}{}
    Soit $u \in \mathscr L(E)$ (resp $A \in \mathscr M_n(\K)$) tel que $\chi_u$ (resp $\chi_A$) soit scindé, 
    c'est-à-dire que $u$ (resp $A$) est trigonalisable.

    Notons $\displaystyle \chi_u = \prod_{i = 1 }^{ p } (X - \lambda_i)^{m_i}$ avec $\lambda_1, \ldots, \lambda_p$ 
    deux à deux distincts. Alors 
    \begin{itemize}
        \item $\tr(u) = \displaystyle \sum_{i = 1 }^{ p } m_i \lambda_i$
        \item $\det(u) = \displaystyle \prod_{i = 1 }^{ p } \lambda_i^{m_i}$
    \end{itemize}
\end{proposition}

\begin{proof}
    Par trigonalisation, il existe $B$ base de $E$ telle que 

    $M_B(u) = \begin{pmatrix}
        \lambda_1 & * & \cdots & * \\
        0 & \lambda_2 & \cdots & * \\
        \vdots & \vdots & \ddots & \vdots \\
        0 & 0 & \cdots & \lambda_n
    \end{pmatrix}$, d'où le résultat par calcul. 
\end{proof}

\begin{corollaire}{LHP}{}
    Soit $A \in \mathscr M_n(\K)$ trigonalisable, alors 

    $A = PTP^{-1}$ avec $T = \begin{pmatrix}
        \lambda_1 & * & \cdots & * \\
        0 & \lambda_2 & \cdots & * \\
        \vdots & \vdots & \ddots & \vdots \\
        0 & 0 & \cdots & \lambda_n
    \end{pmatrix}$ et $P \in GL_n(\K)$
    donc pour $k \in \N$, 
    $$A^k = P T^k P^{-1} \text{ avec } T^k = \begin{pmatrix}
        \lambda_1^k & * & \cdots & * \\
        0 & \lambda_2^k & \cdots & * \\
        \vdots & \vdots & \ddots & \vdots \\
        0 & 0 & \cdots & \lambda_n^k
    \end{pmatrix}$$

    Donc $\tr(A^k) = \displaystyle \sum_{i = 1 }^{ p } m_i\lambda_i^k$.
\end{corollaire}

\begin{remarque}
    Dans les propositions précédentes, les valeurs propres dont l'ordre de multiplicité est 
    supérieur à 1 ont cette valeur propre répétée avec multiplicité sur la 
    diagonale les unes à la suite des autres.
\end{remarque}


\begin{exemple}
    $A = \begin{pmatrix}
        1 & & & & & \\
        & -1 & & & 1 & \\
        & & 1 & & & \\
        & & & \ddots & & \\
        & 1 & & & 1 & \\
        & & & & & -1
        \end{pmatrix} \in \mathscr M_2n(\R)$.

        On cherche les valeurs propres et la dimension des espaces propres. 

        $A + 2I = \begin{pmatrix}
            3 & & & & & \\
            & 1 & & & 1 & \\
            & & 3 & & & \\
            & & & \ddots & & \\
            & 1 & & & 1 & \\
            & & & & & 1
        \end{pmatrix}$ a $n$ colonnes identiques, 
        
        donc est de rang $\leqslant n + 1$.

        Donc (théorme du rang), $\dim \ker(A + 2I) \geqslant 2n - (n + 1) = n - 1$.

        Donc $-2$ valeur propre de multiplicité $\geqslant n - 1$.

        Par le même raisonnement, $A$ possède $n$ colonnes identiques, donc est 
        de rang $\leqslant n + 1$.

        Donc (théorème du rang), $\dim \ker(A - I) \geqslant 2n - (n + 1) = n - 1$.

        Donc $0$ valeur propre de multiplicité $\geqslant n - 1$.

        $A$ est symétrique réelle donc diagonalisable, donc $\chi_A$ est scindé. 

        et $\chi_A = X^{n -1} (X + 2)^{n - 1} (X - a)(X - b)$ avec $a, b \in \R$.

        Or $\tr A = 0 = (n - 1) \cdot 0 + (n - 1)(-2) + a + b$ donc $a + b = 2(n - 1)$.

        $\tr A^2 = (n - 1) 0^2 + (n - 1)(-2)^2 + a^2 + b^2$ 
        or les coefs diagonaux de $A^2$ sont tous égaux à $2n$ donc 
        $(2n)^2 = 4(n - 1) + a^2 + b^2$.

        $\begin{cases}
            a + b = 2(n - 1) \\
            a^2 + b^2 = 4n^2 - 4n + 4
        \end{cases}$ donc $2ab = 4n^2 - 8n + 4 - (4n^2 - 4n + 4)$ et 
        donc $ab = -2n$. 

        Donc $\begin{cases}
            a + b = 2(n - 1) \\
            ab = -2n
        \end{cases}$ donc $a$ et $b$ sont les racines de $P = X^2 - 2(n - 1)X - 2n$.

        $\Delta = 4(n - 1)^2 + 8 n = 4n^2 + 4$ donc 
        $\{a, b\} = \{n - 1 - \sqrt{n^2 + 1}, n - 1 + \sqrt{n^2 + 1}\}$.

        Donc $\chi_A = X^{n -1} (X + 2)^{n - 1} (X - (n - 1 - \sqrt{n^2 + 1}))(X - (n - 1 + \sqrt{n^2 + 1}))$.
\end{exemple}

\section{Polynômes d'endomorphismes et de matrices carrées}
\subsection{Morphisme d'algèbre }

\begin{theoreme}{Morphisme d'algèbre polynôme}{}
    Soit $E$ un $\K$-espace vectoriel, $u \in \mathscr L(E)$ et 
    $P = \displaystyle \sum_{k = 0 }^{ N } a_k X^k \in \K[X]$.

    On notera $P(u) = \displaystyle \sum_{k = 0 }^{ N } a_k u^k \in \mathscr L(E)$.

    Alors l'application $\fonction{\varphi_u}{\K[X]}{\mathscr L(E)}{P}{P(u)}$ est un morphisme d'algèbres.

    Son image est notée $\K[u]$, c'est une sous algèbre commutative de $\mathscr L(E)$
    appelée algèbre des polynômes en $u$.

    Son noyau est un idéal de $\K[X]$ appelé idéal des polynômes annulateurs de $u$.

    \lvspace Soit $A \in \mathscr M_n(\K)$ et $P = \displaystyle \sum_{k = 0 }^{ N } a_k X^k \in \K[X]$.

    On notera $P(A) = \displaystyle \sum_{k = 0 }^{ N } a_k A^k \in \mathscr M_n(\K)$.

    Alors l'application $\fonction{\varphi_A}{\K[X]}{\mathscr M_n(\K)}{P}{P(A)}$ est un morphisme d'algèbres.

    Son image notée $\K[A]$ est une sous algèbre commutative de $\mathscr M_n(\K)$ appelée 
    sous algèbre des polynômes en $A$.

    Son noyau est un idéal de $\K[X]$ appelé idéal des polynômes annulateurs de $A$.

    De plus, si $E$ est de dimension finie $n$, $B$ base de $E$ et $A = M_B(u)$, alors
    $\forall P \in \K[X], \, P(A) = M_B(P(u))$ et donc $\ker \varphi_u = \ker \varphi_A$.
\end{theoreme}

\begin{proof}
    $\varphi_u(1) = \id$ car $\displaystyle 1 = \sum_{k = 0 }^{ 0 } a_k X^k$ avec $a_0 = 1$.

    Soient $P, Q \in \K[X]$, $\displaystyle P = \sum_{k = 0 }^{ N } a_k X^k$ et $\displaystyle Q = \sum_{k = 0 }^{ N} b_k X^k$
    avec $N = \max(\deg P, \deg Q)$ 

    $\alpha, \beta \in \K$ alors $\displaystyle\alpha P + \beta Q = \sum_{k = 0 }^{n } (\alpha a_k + \beta b_k) X^k$.

    Donc 
    \begin{align}
        \varphi_u(\alpha P + \beta Q) & = (\alpha P + \beta Q)(u) \\
        & = \sum_{k = 0 }^{ N } (\alpha a_k + \beta b_k) u^k \\
        & = \alpha \sum_{k = 0 }^{ N } a_k u^k + \beta \sum_{k = 0 }^{ N } b_k u^k \\
        & = \alpha P(u) + \beta Q(u) \\
    \end{align}

    $\displaystyle PQ = \sum_{k = 0 }^{ 2N } c_k X^k$ avec $\displaystyle c_k = \sum_{i = 0 }^{ k } a_i b_{k - i}$.

    Donc $\displaystyle (PQ)(u) = \sum_{k = 0 }^{ 2N } c_k u^k$ c'est à dire 

    $\displaystyle (PQ)(u) = \sum_{k = 0 }^{ N} \sum_{i = 0 }^{k } a_i b_{k - i} u^k$.

    Or $\displaystyle P(u) \circ Q(u) = (\sum_{i = 0 }^{ N } a_i u^i)\circ( \sum_{j = 0 }^{ N } b_j u^j)
    = \sum_{i = 0 }^{ 2N} \sum_{j = 0 }^{2N } a_i b_j u^{i + j}$.
\end{proof}



\begin{exemple}
    Soit $A \in \mathscr M_n(\K)$ telle que $A^2 = 3A - 2 I_n$. 

    Posons $P = X^2 - 3X + 2 = (X - 1)(X - 2)$.

    Alors $P(A) = 0$ et donc $P$ est annulateur de $A$ : $P \in \ker \varphi_A$.

    Effectuons la division euclidienne de $X^k$ par $P$. $X^k = PQ_k + R_k$ avec 
    $\deg R_k < \deg P = 2$.

    Donc $R_k = a_k(X - 1) + b_k (X - 2)$ car $(X - 1)$ et $(X - 2)$ forment une base de $\R_1[X]$.

    Donc $X^k = P Q_k + a_k (X - 1) + b_k (X - 2)$ sont égaux dans $\R[X]$. 

    Appliquons $\varphi_A$ qui est un morphisme d'algèbre : $A^k = P(A) Q_k(A) + a_k (A - I_n) + b_k (A - 2 I_n)$.

    Or $P(A) = 0$ donc $A^k = a_k (A - I_n) + b_k (A - 2 I_n)$.

    Evaluons maintenant la relation en $1$ et $2$ les racines de $P$ : 
    \begin{align}
        1^k & = a_k (1 - 1) + b_k (1 - 2) = -b_k \\
        2^k & = a_k (2 - 1) + b_k (2 - 2) = a_k
    \end{align}

    Donc $A^k = 2^k (A - I_n) - A + 2 I_n$.
\end{exemple}


\subsection{Polynôme minimal }

\begin{theoreme}{Polynôme minimal d'endomorphismes}{}
    \index{Polynôme!minimal}
    Soit $E$ un $\K$-espace vectoriel de dimension finie et $u \in \mathscr L(E)$.\\
    Alors $\ker \varphi_u  = \{P \in \K[X], P(u) = 0\} \neq \{0\}$ donc il existe un unique 
    polynôme unitaire appelé polynôme minimal de $u$ et noté $\mu_u$ tel que 
    $$\ker \varphi_u = \mu_u \K[X]$$
    Cet idéal est appelé idéal des polynômes annulateurs de $u$.
\end{theoreme}

\idee{
    Dimension
}

\begin{proof}
    $\varphi_u: \K[X] \to \mathscr L(E)$ est un morphisme d'algèbres.

    $\im \varphi_u \subset \mathscr L(E)$ donc est de dimension finie.

    Si on avait $\ker \varphi_u = \{0\}$, alors $\K[X]$ en serait un supp, 
    donc (lemme du th du rang), il serait isomorphe à $\im \varphi_u$.

    Or $\K[X]$ de dimension infinie, et $\im \varphi_u$ de dimension finie : absurde. 

    Donc $\ker \varphi_u \neq \{0\}$.

    $\K[X]$ est un anneau principal, donc $\ker \varphi_u$ est engendré par un unique polynôme unitaire
\end{proof}

\begin{theoreme}{Polynôme minimal matriciel}{}
    Soit $E$ un $\K$-espace vectoriel de dimension finie et $A \in \mathscr M_n(\K)$. \\
    Alors $\ker \varphi_A = \{P \in \K[X], P(A) = 0\} \neq \{0\}$ donc il existe un unique
    polynôme unitaire appelé polynôme minimal de $A$ et noté $\mu_A$ tel que
    $$\ker \varphi_A = \mu_A \K[X]$$
    Cet idéal est appelé idéal des polynômes annulateurs de $A$.
\end{theoreme}

\begin{proof}
    Identique à la preuve précédente.
\end{proof}

\begin{theoreme}{Lien polynôme minimal endomorphisme et matrice}{}
    Soit $E$ un $\K$-espace vectoriel de dimension finie, $u \in \mathscr L(E)$, $B$ une base de $E$ et $A = M_B(u)$.
    Alors $\mu_u = \mu_A$.
\end{theoreme}

\begin{proof}
    $\ker \varphi_u = \ker \varphi_A$ donc les polynômes minimaux sont égaux.
\end{proof}

\begin{remarque}
    $\deg \mu_A \geqslant 1$ car sinon $\mu_A$ serait constant égal à $1$ et donc
    $\ker \varphi_A = \K[X]$ ce qui est absurde.

    \avspace Si $\deg \mu_A = 1$ alors $\mu_A = X - \alpha$ et en ce cas 
    $\mu_A(A) = 0$ donc $A = \alpha I_n$.

    Ainsi $\deg \mu_A = 1$ ssi $A$ est la matrice d'une homothétie.

    \avspace Pour trouver $\mu_A$ on peut utiliser l'algo suivant : 
    \begin{itemize}
        \item Si $A$ est ma matrice d'une homothétie, alors $\deg \mu_A = 1$ et $\mu_A = X - \alpha$.
        \item Sinon, si $(\id, A, A^2)$ est liée alors $\deg \mu_A = 2$. Si elle est libre on itère 
            le procédé en considérant $(\id, A, A^2, A^3)$ etc.
    \end{itemize}
\end{remarque}

\begin{exemple}
    $A = \begin{pmatrix}
        1 & 2 \\
        4 & 7
    \end{pmatrix}$

    $A$ n'est pas une homothétie.

    $A^2 = \begin{pmatrix}
        9 & 16 \\
        32 & 57
    \end{pmatrix} = 8 A + I_2$ donc $P = X^2 - 8X - 1$ est annulateur de $A$ et 
    donc $\mu_A = X^2 - 8X - 1$.

    \lvspace $A = \begin{pmatrix}
        1 & 0 & 1 \\
        -1& 2 & 1 \\
        0 & 1 & -1
    \end{pmatrix}$. 

    $A$ n'est pas une homothétie.

    $A^2 = \begin{pmatrix}
        1 & 1 & 0 \\
        -3 & 5 & 1 \\
        -1 & 1 & 2
    \end{pmatrix}$ et $(I_3, A, A^2)$ est libre donc $\deg \mu_A \geqslant 3$.

    $A^3 = \begin{pmatrix}
        0 & 2 & 2 \\
        -8 & 10 & 2 \\
        -2 & 4 & -2
    \end{pmatrix} = 2 A^2 + 2A - 4 I_3$, et donc $P = X^3 - 2X^2 - 2X + 4$ est annulateur de $A$.

    De plus il est de degré minimal, donc $\mu_A = X^3 - 2X^2 - 2X + 4$.
\end{exemple}

\begin{theorement}{dimension de l'algèbre des polynômes d'un endomorphisme}{}
    Soit $u \in \mathscr L(E)$ un $\K$-espace vectoriel de dimension finie, 
    de polynôme minimal $\mu_u$. 

    Alors $\dim \K[u] = \deg \mu_u$ et $(\id, u, \ldots, u^{\deg \mu_u - 1})$ est une base de $\K[u]$.
\end{theorement}

\begin{theorement}{dimension de l'algèbre des polynômes d'une matrice}{}
    Soit $A \in \mathscr M_n(\K)$ de polynôme minimal $\mu_A$.

    Alors $\dim \K[A] = \deg \mu_A$ et $(I_n, A, \ldots, A^{\deg \mu_A - 1})$ est une base de $\K[A]$.
\end{theorement}

\begin{proof}
    Considérons $d = \deg \mu_A$ et $\fonction{\tilde \varphi_A}{\K_{d - 1}[X]}{\mathbb K[A]}{P}{P(A)}$.

    $\tilde \varphi_A$ est une application linéaire. 

    Si $P \in \ker \tilde \varphi_A$, alors $P(A) = 0$ et donc $\mu_A | P$.

    Or $\deg \mu_A > \deg P$ et donc $P = 0$.

    Donc $\tilde \varphi_A$ est injective.

    \avspace Soit $B \in \K [A]$, alors $\exists P \in \K[X]$ tel que $B = P(A)$.

    Par division euclidienne, $P = \mu_A Q + R$ avec $\deg R \leqslant d - 1$, donc 
    (par morphisme d'algèbre) $P(A) = R(A)$.

    Donc $B = R(A) = \tilde \varphi(R)$ donc $B \in \im \tilde \varphi_A$ et donc 
    $\tilde \varphi_A$ est surjective.
\end{proof}

\begin{exemple}
    Si $p$ est un projecteur de $E$, alors $p^2 = p$ donc $X^2 - X$ est annulateur de $p$.

    Donc $\mu_P \in \{X^2 - X, X - 1, X\}$. 3 cas : 
    \begin{itemize}
        \item $p = 0$
        \item $p = \id$
        \item $\mu_P$ = $X^2 - X$ 
    \end{itemize}

    Dans le cas où $p \notin \{0, \id\}$, on a $\K[p]$ de dimension $2$ qui a pour base $(\id, p)$.

    \lvspace Soit $A \in \mathscr M_2(\K)$ telle que $A = \begin{pmatrix}
        a & b \\
        c & d
    \end{pmatrix}$ et donc $A^2 = \begin{pmatrix}
        a^2 + bc & ab + bd \\
        ac + cd & bc + d^2
    \end{pmatrix} = (a + d) A  + (b c - ad) I_2 = \tr(A) A + \det(A) I_2$.

    Donc $X^2 - \tr(A) X + \det(A)$ est annulateur de $A$.

    $2$ cas possibles : 
    \begin{itemize}
        \item $A$ est la matrice d'une homothétie, $A = \lambda I_2$ et $\mu_A = X - \lambda$ donc 
        $\K[A] = \Vect(I_2)$ de dimension $1$.
        \item $A$ n'est pas la matrice d'une homothétie, donc $\mu_A = X^2 - \tr(A) X + \det(A)$ et 
        $\K[A] = \Vect(I_2, A)$ de dimension $2$.
    \end{itemize}
\end{exemple}

\subsection{Lien avec le polynôme caractéristique}

\begin{proposition}{Polynôme minimal et polynôme caractéristique}{}
    Soient $E$ un $\K$-espace vectoriel, $u \in \mathscr L(E)$, $x \in E$ et 
    $\lambda \in \K$. Supposons que $u(x) = \lambda x$. Alors $(P(u))(x) = P(\lambda)\cdot x$.
\end{proposition}

\begin{proof}
    $\forall n \in \N, \, X^n(u)(x) = u^n(x) = \lambda^n x$.

    Par récurrence sur $n$ : $n = 0$ : $\id(x) = x = \lambda^0 x$.

    HR : Supposons que $u^n(x) = \lambda^n x$. Alors $u^{n + 1}(x) = u(u^n(x)) = u(\lambda^n x) = \lambda^n u(x) = \lambda^{n + 1} x$.

    Soit $P = \displaystyle \sum_{k = 0 }^{ N } a_k X^k \in \K[X]$.

    Alors $(P(u))(x) = \sum_{k = 0 }^{ N } a_k u^k(x) = \sum_{k = 0 }^{ N } a_k \lambda^k x = P(\lambda) \cdot x$.
\end{proof}

\begin{proposition}{Polynôme minimal et valeurs propres}{}
    Soit $P \in \K[X], \, u \in \mathscr L(E)$. Si $P$ annule $u$, alors toute valeur propre de 
    $u$ est racine de $P$.
\end{proposition}

\begin{proof}
    Supposons que $P(u) = 0$, et soit $\lambda $ une valeur propre de $u$. Alors 
    $\exists x \in E$ tel que $x \neq 0$ et $u(x) = \lambda x$.

    Donc $(P(u))(x) = P(\lambda) \cdot x$.

    Donc $P(\lambda)x = 0$, or $x \neq 0$ donc $P(\lambda) = 0$.
\end{proof}

\begin{exemple}
    Si $p$ est un projecteur non trivial, alors $P = X^2 -X$ annule $p$. Les valeurs propres de 
    $p$ qui sont $0$ et $1$ sont bien racines de $P$.
\end{exemple}

\begin{proposition}{Polynôme minimal et valeurs propres}{}
    Soit $E$ un $\K$-espace vectoriel de dimension finie, et $u \in \mathscr L(E)$ de polynôme minimal $\mu_u$.

    Alors les racines de $\mu_u$ dans $\K$ sont les valeurs propres de $u$.
\end{proposition}

\begin{proof}
    $\mu_u$ annule $u$ donc les valeurs propres de $u$ sont racines de $\mu_u$ (par proposition précédente).

    Soit $\lambda \in \K$ tel que $\lambda$ n'est pas valeur propre de $u$. 

    Alors $\det(u - \lambda \id) \neq 0$ donc $u - \lambda \id$ est inversible.

    Si on avait $\lambda $ racine de $\mu_u$ on aurait $\mu_u = (X - \lambda)^k \cdot P$ avec $k \geqslant 1$ et 
    $\deg P < \deg \mu_u$

    Donc $(u - \lambda \id)^k \circ P(u) = 0$. 

    Donc $P(u) = 0$, absurde car $\mu_u$ est de degré minimal.

    Donc $\lambda$ n'est pas racine de $\mu_u$.
\end{proof}

\begin{theoreme}{Cayley Hamilton}{}
    Soit $E$ un $\K$-espace vectoriel de dimension finie et $u \in \mathscr L(E)$.

    Alors $\mu_u | \chi_u$, c'est à dire que $\chi_u(u) = 0$. 

    Le polynôme caractéristique est un polynôme annulateur de $u$, et de même 
    pour une matrice $A \in \mathscr M_n(\K)$, $\chi_A(A) = 0$.
\end{theoreme}

\idee{Construire une base avec des vecteurs de la forme $x, u(x), \ldots, u^k(x)$.}

\begin{proof}
    Non exigible.

    Soit $u \in \mathscr L(E)$, $E$ de dimension finie. On veur montrer que 
    $\forall x \in E, \, (\chi_u(u))(x) = 0$.

    Si $x = 0$, c'est vrai. 

    Sinon considérons les vecteurs $e_0 = x$, $e_1 = u(x)$, $\ldots$, $e_k = u^k(x)$ 
    et $A = \{k \in \N \mid (x, u(x) \cdots, u^k(x)) \text{ est libre }\}$

    $A \subset \N, \, A \neq \emptyset$ car $0 \in A$.

    $A$ est majorée car $E$ de dimension finie. Posons $N = \max A$, alors $(e_0, e_1, \ldots, e_N)$ est libre et
    $(e_0, e_1, \ldots, e_{N + 1})$ est liée.

    Donc $\exists \alpha_0, \ldots, \alpha_N \in \K$ tels que
    $e_{N + 1} = \sum_{k = 0 }^{ N } \alpha_k e_k$.

    Posons $F = \Vect(e_0, e_1, \ldots, e_N)$ espace de dimension $N + 1$. 

    $\forall i \in \ninter 0 N$, $u(e_i) \in F$ donc $F$ stable par $u$.

    Considérons l'endomorphisme $\tilde u $ induit par $u$ sur $F$.

    $M_B(u) = \begin{pmatrix}
        0 & 0 & \cdots & 0 & \alpha_0 \\
        1 & 0 & \cdots & 0 & \alpha_1 \\
        0 & 1 & \cdots & 0 & \alpha_2 \\
        \vdots & \vdots & \ddots & \vdots & \vdots \\
        0 & 0 & \cdots & 1 & \alpha_N
    \end{pmatrix}$ matrice compagnon de 
    
    $\displaystyle Q(X) = X^{N + 1} - \sum_{k = 0 }^{ N } \alpha_k X^k$.

    Or $\chi_{\tilde u} \mid \chi_u$ donc $\exists P \in \K[X], \, \chi_u = P \cdot \chi_{\tilde u}$.

    Donc 

    \begin{align}
        (\chi_u(u))(x) & = ((P \cdot \chi_{\tilde u})(u))(x) \\
        & = (P(u) \circ \chi_{\tilde u}(u))(x) \\
        & = P(u)(\chi_{\tilde u}(u)(x)) \\
    \end{align}

    Or $\displaystyle (\chi_{\tilde u}(\tilde u))(x) = (u^{N + 1} - \sum_{k = 0 }^{ N } \alpha_k u^k)(x) = e_{N + 1} - \sum_{k = 0 }^{ N } \alpha_k e_k = 0$.
\end{proof}

\begin{corollaire}{polynôme minimal et polynôme caractéristique}{}
    Soit $u \in \mathscr L(E)$ un $\K$-espace vectoriel de dimension finie $n$.

    Dans le cas où $\chi_u$ est scindé (ce qui est toujours le cas si $\K = \C$), 
    notons $\Sp(u) = \{\lambda_1, \ldots, \lambda_p\}$. On a alors : 
    $$\chi_u = \displaystyle \prod_{j = 1 }^{ p } (X - \lambda_j)^{m_j}$$

    Avec $\ds \sum_{j = 1 }^{p } m_j = n$.
    Alors : 
    $$\mu_u = \displaystyle \prod_{j = 1 }^{ p } (X - \lambda_j)^{k_j}$$

    De plus on a : 
    \begin{itemize}
        \item $\forall j \in \ninter{1}{p}, \, 1 \leqslant k_j \leqslant m_j$
        \item $1 \leqslant \dim E_{\lambda_j}(u) \leqslant m_j$ 
    \end{itemize}
\end{corollaire}

\subsection{Lien avec la réduction}

\begin{theoreme}{Lemme de décomposition des noyaux}{}
    Soient $P_1, \dots, P_k \in \K[X]$ deux à deux premiers entre eux et 
    $$P = \displaystyle \prod_{j = 1 }^{ k } P_j$$   
    
    Soit $u \in \mathscr L(E)$ alors 
    $$\displaystyle \ker P(u) = \bigoplus_{j = 1 }^{ k } \ker P_j(u)$$
\end{theoreme}

\idee{Bézout}

\begin{proof}
    Par récurrence sur $k$. 

    Si $k = 1$, sans intérêt. 

    Si $k = 2$, soit $P_1, P_2 \in \K[X]$ premiers entre eux et $P = P_1 P_2$.

    Soit $x \in \ker P(u)$. 

    Par le théorème de Bézout, il existe $U, V \in \K[X]$ tels que $U P_1 + V P_2 = 1$.

    Par le morphisme d'algèbre $P \mapsto P(u)$, on a 
    $\id_E = P_1(u) \circ U(u) + P_2(u) \circ V(u)$ donc 

    $$x = \underbrace{(P_1(u))((U(u))(x))}_{= x_1} + \underbrace{(P_2(u))((V(u))(x))}_{= x_2}$$

    \begin{align}
        (P_1(u))(x_2) & = (P_1(u) \circ P_2(u) \circ V(u))(x) \\
        & = ((P_1 P_2V)(u))(x) \\
        & = ((PV)(u))(x) \\
        & = 0
    \end{align}

    Donc $x_2 \in \ker P_1(u)$. De même, $x_1 \in \ker P_2(u)$.

    Donc $\ker P(u) \subset \ker P_1(u) + \ker P_2(u)$.


    \lvspace Soit $x \in \ker P_1(u) + \ker P_2(u)$, alors $\exists y \in \ker P_1 (u), 
    \exists z \in \ker P_2(u)$ tels que $x = y + z$.

    $P(u)(x) = (P(u))(y) + (P(u))(z) = P_2(u)(P_1(u)(y)) + P_1(u)(P_2(u)(z)) = 0$. 

    Donc $\ker P_1 (u) + \ker P_2(u) \subset \ker P(u)$ d'où l'égalité. 

    Soit $x \in \ker P_1(u) \cap \ker P_2(u)$. Par bézout, 

    $x = U(u)(P_1(u)(x)) + V(u)(P_2(u)(x)) = 0$ donc l'intersection est réduite à $\{0\}$.

    \uindent{Hypthèse de récurrence } Supposons le résultat vrai pour $k$ polynômes. 

    Soient $P_1, \dots, P_{k + 1} \in \K[X]$ deux à deux premiers entre eux et posons 
    $P = \displaystyle \prod_{i = 1}^{ k + 1 } P_i$ et $Q = \displaystyle \prod_{i = 1 }^{ k } P_i$.

    Alors $P = Q \cdot P_{k + 1}$ et $Q \wedge P_{k + 1} = 1$ avec le résultat pour $k = 2$.

    $\ker P(u) = \ker Q(u) \oplus \ker P_{k + 1}(u)$. Or (HR), 
    $\ker(Q(u)) = \bigoplus_{i = 1 }^{ k } \ker P_i(u)$ et la conclusion par associativité de 
    la somme directe.
\end{proof}


\begin{corollaire}{}{}
    Soit $u \in \mathscr L(E)$ et $P$ polynôme annulateur de $u$. Supposons que 
    $P = \displaystyle \prod_{i = 1}^{ k } P_i$ avec $P_1, \ldots, P_k$ deux à deux premiers entre eux.

    Alors par le lemme des noyaux, $\ker P(u) = \bigoplus_{i = 1 }^{ k } \ker P_i(u)$.

    Or $P(u) = 0$ donc $\ker P(u) = E$. Donc 
    $$E = \bigoplus_{i = 1 }^{ k } \ker P_i(u)$$
\end{corollaire}

\begin{theoreme}{Caractérisation de la diagonalisabilité}{}
    Soit $E$ un $\K$-espace vectoriel de dimension finie et $u \in \mathscr L(E)$. 
    Alors les assertions suivantes sont équivalentes :
    \begin{enumerate}
        \item $u$ est diagonalisable
        \item Le polynôme minimal de $u$ est scindé à racines simples.
        \item $u$ annule un polynôme scindé à racines simples.
    \end{enumerate}
\end{theoreme}

\begin{proof}
    $(2) \implies (3)$ : car $\mu_u$ annule $u$.

    \lvspace Montrons que $3 \implies 2$.
    
    Supposons (3). Soit $P$ un polynôme scindé à racines simples annulant $u$. Alors $\mu_u \mid P$ donc 
    $\mu_u$ est scindé à racines simples et donc (2).

    \lvspace Montrons que $1 \implies 3$. 
    
    Supposons (1). 
    Alors $E = \displaystyle \bigoplus_{i = 1 }^{ p } E_{\lambda_i}(u)$ avec $\Sp(u) = \{\lambda_1, \ldots, \lambda_p\}$.

    Posons $P = \displaystyle \prod_{i = 1 }^{ p } (X - \lambda_i)$.

    Soit $x \in E$ décomposé en $x = \displaystyle \sum_{i = 1 }^{ p } x_i$ avec $x_i \in E_{\lambda_i}(u)$.

    Alors $P(u)(x) = \displaystyle \sum_{i = 1 }^{ p } P(u)(x_i)$, or $P(u)(x_i) = (Q_i(u) \circ (u - \lambda_i \id))(x_i)$
    en notant $P = Q_i \times (X - \lambda_i)$.

    Donc $P(u)(x_i) = 0$, donc $P(u) = 0$ et donc (3).

    \lvspace Montrons que $3 \implies 1$. 
    
    Supposons (3). 

    Soit $P$ scindé à racines simples telles que $P = \displaystyle \prod_{i = 1 }^{ k } (X - \lambda_i)$
    et $P(u) = 0$.

    Les $(X - \lambda_i)$ sont des polynômes deux à deux premiers entre eux. Donc 
    (lemme des noyaux) $E = \displaystyle \bigoplus_{i = 1 }^{ k } \ker (u - \lambda_i \id)$.

    Notons $\lambda_1, \dots, \lambda_p$ les $\lambda_i$ pour lesquels ce noyau est non trivial.

    Donc $E = \displaystyle \bigoplus_{i = 1 }^{ p } \ker (u - \lambda_i \id) $ d'où le (1).
\end{proof}

\begin{exemple}
    $A = \begin{pmatrix}
        1 & 2 & 4 \\
        -1 & 3 & 5 \\
        3 & 1 & 3 
    \end{pmatrix}$ eet $L_3 + L_2 = 2L_1$. $A$ n'est pas inversible, donc $0$ est valeur propre de $A$.

    $\begin{pmatrix}
        1 \\ 1 \\ 1
    \end{pmatrix}$ est valeur propre pour $7$. 

    Donc $\chi_A(X) = X(X - 7)(X - \alpha)$ et $0 + 7 + \alpha = \tr(A) = 7$ donc $\alpha = 0$.

    Donc $\chi_A(X) = X^2 (X - 7)$ et $\dim E_7(A) = 1$ car valeur propre simple.

    $\dim \ker A \geqslant 1$ et $\rg A \leqslant 2$, or $\rg A \notin \{0, 1\}$ (car $L_1$ et $L_2$
    sont libres) donc $\rg A = 2$ et donc $\dim \ker A = 1$.

    $A$ non diagonalisable or $\mu_A \mid \chi_A$ et a les mêmes racines et unitaire donc 
    $\mu_A \in \{X(X - 7), X^2(X - 7)\}$ or $A$ non diagonalisable sonc $\mu_A$ n'est pas scindé à racines simples donc 
    $\mu_A = X^2 (X - 7)$.
\end{exemple}

\begin{proposition}{Polynôme minimal et sous-espace stable}{}
    Soit $u \in \mathscr L(E)$ avec $E$ un $\K$-espace vectoriel de dimension finie. 

    Soit $F$ sous-espace vectoriel de $E$ stable par $u$, et $\tilde u$ l'endomorphisme induit. 

    Alors $\mu_{\tilde u} \mid \mu_u$.
\end{proposition}


\begin{proof}
    Soit $x \in F$.

    $(\mu_u(\tilde u))(x) = (\mu_u(u))(x) = 0$ donc $\mu_u(\tilde u) = 0$.

    Donc $\mu_{\tilde u} \mid \mu_u$.
\end{proof}

\begin{corollaire}{endomorphisme induit et diagonalisabilité}{}
    Un endomorphisme induit par un endomorphisme diagonalisable est lui aussi diagonalisable.
\end{corollaire}

\begin{proof}
    Avec les notations précédentes : $u$ est diagonalisable donc $\mu_u$
    scindé à racines simples, or $\mu_{\tilde u} \mid \mu_u$ donc $\mu_{\tilde u}$ 
    scindé à racines simples donc $\tilde u$ est diagonalisable.
\end{proof}

\begin{corollaire}{LHP}{}
    \index{Codiagonalisable}
    Soient $u, v \in \mathscr L(E)$ tels que $u \circ v = v \circ u$ et 
    $u$ et $v$ diagonalisables. Alors il existe $B$ une base de $E$ telle que 
    $M_B(u)$ et $M_B(v)$ sont diagonales.

    On dit que $u$ et $v$ sont codiagonalisables.
\end{corollaire}

\begin{proof}
    Notons $\lambda_1, \dots, \lambda_p$ les valeurs propres de $u$. 

    $u \circ v = v \circ u$ donc $\forall j \in \ninter{1}{p}, \, E_{\lambda_j}(u)$ est stable par $v$.

    Notons $v_j$ l'endomorphisme induit par $v$ sur $E_{\lambda_j}(u)$.

    $v$ étant diagonalisable, $v_j$ l'est aussi. 

    Notons $B_j$ une base de $E_{\lambda_j}(u)$ formée de vecteurs propres de $v_j$.

    Les vecteurs de $B_j$ sont propres pour $u$ et $v$, or $u$ diagonalisable donc 
    $E = \displaystyle \bigoplus_{j = 1 }^{ p } E_{\lambda_j}(u)$, donc 
    $B = B_1 \cup \dots \cup B_p$ base de $E$ donc $M_B(u)$ et $M_B(v)$ sont diagonales.
\end{proof}

\begin{theoreme}{Caractérisation de la trigonalisabilité}{}
    Soit $u \in \mathscr L(E)$, alors les assertions suivantes sont équivalentes :  
    \begin{enumerate}
        \item $u$ est trigonalisable
        \item $\mu_u$ est scindé
        \item $u$ annule un polynôme scindé
    \end{enumerate}
\end{theoreme}

\begin{proof}
    Montrons que $1 \implies 3$.
    
    Si $u$ est trigonalisable alors $\chi_u$ scindé, or $\chi_u$ annule $u$ d'où (3).

    \lvspace Montrons que $3 \implies 2$.
    
    Supposons (3). Soit $P$ scindé tel que $P(u) = 0$, or $\mu_u \mid P$ donc $\mu_u$ est scindé d'où (2).

    \lvspace Montrons que $2 \implies 1$.
    
    Supposons le (2). $\mu_u$ est scindé, donc $\mu_u$ a les mêmes racines dans $\K$ et $\C$, 
    $\mu_u$ et $\chi_u$ ont les mêmes racines, donc $\chi_u$ a les mêmes racines dans $\K$ et $\C$.

    Donc $\chi_u$ scindé. Donc $u$ trigonalisable. 
\end{proof}

\subsection{Nilpotence}

\begin{definition}{Endomorphisme et matrice nilpotent}{}
    \index{Endomorphisme!nilpotent}
    \index{Matrice!nilpotente}
    Soit $E$ un $\K$-espace vectoriel de dimension finie et $u \in \mathscr L(E)$. 

    On dit que $u$ est nilpotent s'il existe $k \geqslant 1$ tel que $u^k = 0$.

    Le plus petit entier vérifiant $u^k = 0$ est appelé indice de nilpotence de $u$.
    
    De même pour les matrices. 
\end{definition}


\begin{proposition}{Polynôme minimal et nilpotence}{}
    Soit $u \in \mathscr L(E)$, alors $u$ est nilpotent d'ordre $k$ ssi $\mu_u = X^k$.
    De même pour les matrices.
\end{proposition}

\begin{proof}
    Si $u$ nilpotent d'ordre $k$ alors $u^k = 0$, donc $X^k$ annule $u$, 
    donc $\mu_u \mid X^k$, donc $\mu_u \in \{X, X^2, \ldots, X^k\}$.

    Or $u^{k - 1} \neq 0$ donc $\mu_u \notin \{X, X^2, \ldots, X^{k - 1}\}$, donc $\mu_u = X^k$.

    Réciproquement, si $\mu_u = X^k$ alors $u^k = 0$ et $u^{k - 1} \neq 0$ donc $u$ est nilpotent d'ordre $k$.
\end{proof}

\begin{corollaire}{indice de nilpotence et dimension}{}
    Soit $u \in \mathscr L(E)$ nilpotent d'indice $k$ alors $k \leqslant n$.
\end{corollaire}

\begin{proof}
    $\mu_u = X^k \mid \chi_u$ or $\deg \chi_u = n$ donc $k \leqslant n$. 
\end{proof}

\begin{corollaire}{caractérisation de la nilpotence}{}
    Soit $u \in \mathscr L(E)$, alors $u$ est nilpotent ssi $u$ est trigonalisable et $0$ est son unique valeur
    propre. 
\end{corollaire}

\begin{proof}
    Si $u$ est nilpotent, alors $\exists k, \, \mu_u = X^k$ qui est scindé, donc $u$ est trigonalisable.
    De plus les valeurs propres de $u$ sont les racines de $\mu_u$, donc $0$ est l'unique valeur propre.

    Réciproquement, si $u$ est trigonalisable et $0$ unique valeur propre, alors $\mu_u$ est scindé avec 
    $0$ comme unique racine. 

    Donc $\exists k, \, \mu_u = X^k$ donc $u$ est nilpotent.
\end{proof}

\begin{corollaire}{}{}
    Soit $u \in \mathscr L(E)$, $n = \dim E$. $u$ est nilpotent ssi $\chi_u = X^n$.
\end{corollaire}

\begin{proof}
    Si $u$ est nilpotent, $u$ trigonalisable et $\Sp(u) = \{0\}$ donc $\chi_u$ scindé avec $0$ pour unique 
    racine donc $\chi_u = X^n$.

    Si $\chi_u = X^n$ comme $\mu_u \mid \chi_u$, $\exists k, \, \mu_u = X^k$ donc $u$ nilpotente. 
\end{proof}

\begin{remarque}
    Tous ces résultats s'appliquent aux matrices. 
\end{remarque}

\begin{exemple}
    $\begin{pmatrix}
        -2 & 4 \\ 
        -1 & 2
    \end{pmatrix}$ et $\chi_A = \begin{vmatrix}
        X + 2 & -4 \\
        1 & X - 2
    \end{vmatrix} = X^2$ donc $\mu_A \in \{X, X^2\}$.

    $A \neq 0$ donc $\mu_A = X^2$ et $A$ est nilpotente d'indice $2$.

    Donc $A$ non diagonalisable (car la seule matrice nilpotente diagonalisable 
    est la matrice nulle dont le pynôme minimal est $X$).

    $\chi_A$ est scindé, donc $A$ est trigonalisable. 

    $\begin{pmatrix}
        x \\ y
    \end{pmatrix} \in \ker A$ ssi $-2x + 4y = 0$

    $\ker A = \Vect e_1$ avec $e_1 = \begin{pmatrix}
        2 \\ 1
    \end{pmatrix}$.

    Posons $e_2 = A \begin{pmatrix}
        -1 \\ 0
    \end{pmatrix}$ et alors $A = PTP^{-1}$ avec $P = (e_1, e_2)$ et 
    $T = \begin{pmatrix}
        0 & 1 \\
        0 & 0
    \end{pmatrix}$.
\end{exemple}


\subsection{Sous-espaces caractéristiques}

Dans toute cette partie, on suppose que le polynôme caractéristique est scindé.

\begin{definition}{Sous-espace caractéristique}{}
    \index{Sous-espace!caractéristique}
    Soit $u \in \mathscr L(E)$, $\K$-espace vectoriel de dimension finie. Supposons
    que $\chi_u$ est scindé, $\displaystyle \chi_u = \prod_{i = 1 }^{ p } (X - \lambda_i)^{m_i}$.

    \marginenote{Remarque}{$\ker (u - \lambda_i \id )^{m_i}$ désigne le noyau d'une composition, et non pas d'un produit.}
    \lvspace \avspace 

    Pour $i \in \ninter{1}{p}$, on définit le sous-espace caractéristique associé à $\lambda_i$ par 
    $$C_{\lambda_i} = \ker (u - \lambda_i \id )^{m_i}$$

    On a en particulier $E_{\lambda_i}(u) \subset C_{\lambda_i} $
\end{definition}

\begin{theoreme}{Somme directe des sous-espaces caractéristiques}{}
    Soit $u \in \mathscr L(E)$ tel que $\chi_u$ est scindé. Notons 
    $\chi_u = \displaystyle \prod_{i = 1 }^{ p } (X - \lambda_i)^{m_i}$ avec 
    $\lambda_1, \ldots, \lambda_p$ valeurs propres distinctes de $u$
    et pour tout $i$, $C_{\lambda_i}(u) = \ker (u - \lambda_i \id)^{m_i}$.
    Alors $$E = \displaystyle \bigoplus_{i = 1 }^{ p } C_{\lambda_i}(u)$$

    L'espace est donc somme directe des sous-espaces caractéristiques.
\end{theoreme}

\idee{
    On applique le lemme des noyaux au polynôme caractéristique.
}

\begin{proof}
    Posons pour $i \in \ninter 1 k $, $P_i = (X - \lambda_i)^{m_i}$, et
    pour $i \neq j$, $\lambda_i \neq \lambda_j$ donc $P_i$ et $P_j$ sont premiers entre eux.

    Donc par le lemme des noyaux, $\displaystyle \ker \left(\prod_{i = 1 }^{ p } P_i(u) \right) = \bigoplus_{i = 1 }^{ p } \ker P_i(u)$.
    
    Donc $\ds \ker \chi_u(u) = \bigoplus_{i = 1 }^{ p } C_{\lambda_i}(u)$, or $\chi_u$ annule $u$ donc 
    $\ker \chi_u(u) = E$ d'où le résultat.

\end{proof}

\begin{corollaire}{dimension des sous-espaces caractéristiques}{}
    Avec les notations du théorème précédent, $$\forall i \in \ninter{1}{p}, 
    \dim C_{\lambda_i}(u) = m_i$$
\end{corollaire}

\begin{proof}
    On a $E = \displaystyle \bigoplus_{i = 1 }^{ p } C_{\lambda_i}(u)$ qui sont 
    des sous-espaces stables. 

    Dans une base adaptée à cette somme directe, en notant pour tout $i$ 
    $\tilde u_i$ l'endomorphisme induite par $u$ sur $C_{\lambda_i}(u)$, la 
    matrice de $u$ est de la forme

    $M_B(u) = \begin{pmatrix}
        M_1 & 0 & \cdots & 0 \\
        0 & M_2 & \cdots & 0 \\
        \vdots & \vdots & \ddots & \vdots \\
        0 & 0 & \cdots & M_p
    \end{pmatrix}$ avec $M_i = M_{\mathcal B_i}(\tilde u_i)$.

    Donc $M_B(u)$ est diagonale par blocs, et donc 

    \begin{align}
        \chi_u & = \det (X I_n - M_B(u)) \\
        & = \prod_{i = 1 }^{ p } \det (X I_{\dim C_{\lambda_i}} - M_{B_i}(\tilde u_i)) \\
        & = \prod_{i = 1 }^{ p } \chi_{\tilde u_i}
    \end{align}

    On a pour $i \in \ninter 1 p$ et $x \in C_{\lambda_i}(u)$,
    $(\tilde u_i - \lambda_i \id)^{m_i}(x) = 0$ et notons $n_i = \tilde u_i - \lambda_i \id_{C_{\lambda_i}(u)} \in \mathscr L(C_{\lambda_i}(u))$.

    On a alors $n_i^{m_i} = 0$ donc $n_i$ est nilpotent, et on a 
    $P = (X - \lambda_i)^{m_i}$ annulateur de $\tilde u_i$.

    Donc le polynôme minimal de $\tilde u_i$ est de la forme $(X - \lambda_i)^{\alpha_i}$.

    Donc $\lambda_i$ est l'unique valeur propre de $\tilde u_i$, donc $\chi_{\tilde u_i}$ 
    est scindé car $\chi_u$ l'est. 

    Donc $\chi_{\tilde u_i} = (X - \lambda_i)^{\dim C_{\lambda_i}(u)}$, donc 
    $\displaystyle \chi_u = \prod_{i = 1 }^{ p } (X - \lambda_i)^{\dim C_{\lambda_i}(u)}$, 
    et donc par unicité de la décomposition en facteurs premiers,
    $\forall i \in \ninter{1}{p}, \, \dim C_{\lambda_i}(u) = m_i$.
\end{proof}

\begin{corollaire}{sous-espaces caractéristiques et diagonalisabilité}{}
    Soit $u \in \mathscr L(E)$ diagonalisable et $\Sp (u) = \{\lambda_1, \ldots, \lambda_p\}$
    alors 
    
    $$\forall i \in \ninter{1}{p}, \, E_{\lambda_i}(u) = C_{\lambda_i}(u)$$
\end{corollaire}

\begin{proof}
    Comme $u$ est diagonalisable, $\chi_u$ est scindé et 
    
    $\forall i \in \ninter{1}{p}, \, \dim E_{\lambda_i}(u) = m_i$.

    Or $\dim C_{\lambda_i}(u) = m_i$ et $E_{\lambda_i}(u) \subset C_{\lambda_i}(u)$ donc 
    $E_{\lambda_i}(u) = C_{\lambda_i}(u)$.
\end{proof}


\subsubsection{Matrices diagonales par blocs}

Voir quoi faire de ce pavé, si c'est une proposition ou pas et si le titre de la partie a un sens 

Soit $A \in \mathscr M_n(\K)$ associée à $u \in \mathscr L(E)$ dans une base. On suppose 
$\chi_A$ scindé.

Notons $\Sp (A) = \Sp(u) = \{\lambda_1, \ldots, \lambda_p\}$ et 
$\chi_A = \chi_u = \displaystyle \prod_{i = 1 }^{ p } (X - \lambda_i)^{m_i}$.

Notons pour tout $i$ : 
\begin{itemize}
    \item $C_{\lambda_i}(u) = \ker (u - \lambda_i \id)^{m_i}$
    \item $\tilde u_i$ l'endomorphisme induit par $u$ sur $C_{\lambda_i}(u)$
    \item $n_i = \tilde u_i - \lambda_i \id_{C_{\lambda_i}(u)}$
\end{itemize}

Pour tout $i$, $n_i$ est nilpotent d'indice $\leqslant m_i$, 
$n_i$ est donc trigonalisable (car de polynôme caractéristique $(X)^{m_i}$ scindé).

Considérons $B_i$ base de $C_{\lambda_i}(u)$ dans laquelle la matrice de $n_i$ est 
$$N_i = \begin{pmatrix}
    0 & * & \cdots & * \\
    0 & 0 & \cdots & * \\
    \vdots & \vdots & \ddots & \vdots \\
    0 & 0 & \cdots & 0
\end{pmatrix}$$ matrice triangulaire stricte. 

La matrice de $\tilde u_i$ dans cette base est donc 

$$M_i = \begin{pmatrix}
    \lambda_i & * & \cdots & * \\
    0 & \lambda_i & \cdots & * \\
    \vdots & \vdots & \ddots & \vdots \\
    0 & 0 & \cdots & \lambda_i
\end{pmatrix}$$ matrice triangulaire de diagonale $\lambda_i$.

Donc la matrice de $u$ dans la base $B = B_1 \cup \cdots \cup B_p$ est 

$$A' = M_B(u) = \begin{pmatrix}
    M_1 & 0 & \cdots & 0 \\
    0 & M_2 & \cdots & 0 \\
    \vdots & \vdots & \ddots & \vdots \\
    0 & 0 & \cdots & M_p
\end{pmatrix}$$ matrice triangulaire de diagonale constituée des valeurs propres de $u$
et donc $A$ est semblable à $A'$. 

\begin{exemple}
    $A = \begin{pmatrix}
        3 & 9 & 3 \\
        -3 & 6 & 12 \\
        1 & 8 & 6
    \end{pmatrix}$.

    On a $A \cdot \begin{pmatrix}
        1 \\ 1 \\ 1
    \end{pmatrix} = 15 \cdot \begin{pmatrix}
        1 \\ 1 \\ 1
    \end{pmatrix}$ donc $15$ est valeur propre.

    De plus $3 L_3 = 2 L_1 + L_2$ donc $\det(A) = 0$ et $0$ est valeur propre.

    $\chi_A = X(X - 15)(X - \alpha)$ et $\tr(A) = 15$ donc $\alpha = 0$.

    Donc $\chi_A = X^2 (X - 15)$. 

    $0$ est valeur propre double mais $\dim \ker A = 3 - \rg A = 1$ donc non diagonalisable.

    $\ker A = E_0(A) \subsetneq C_0(A) = \ker (A^2)$.

    $E_{15}(A) = C_{15}(A)$ de dimension 1 et base $\begin{pmatrix}
        1 \\ 1 \\ 1
    \end{pmatrix}$.

    $\begin{pmatrix}
        x \\ y \\ z
    \end{pmatrix} \in E_0(A)$ ssi $\begin{cases}
        3x + 9y + 3z = 0 \\
        -3x + 6y + 12z = 0 \\
        x + 8y + 6z = 0
    \end{cases}$ ssi $\begin{cases}
        y + z = 0 \\
        x + 2y = 0
    \end{cases}$

    $E_0(A) = \Vect \begin{pmatrix}
        2 \\ -1 \\ 1
    \end{pmatrix}$.

    Posons $e_1 = \begin{pmatrix}
        1 \\ 1 \\ 1
    \end{pmatrix}$, $e_2 = \begin{pmatrix}
        2 \\ -1 \\ 1
    \end{pmatrix}$ et cherchons $e_3$ dans $\ker A^2$ non colinéaire à $e_2$.

    \avspace Solution 1 : Calculer $A^2$ 

    \avspace Solution 2 : remarquons que $A(e_3) \in \ker A$ donc $A(e_3) $ est colinéaire à $e_2$.

    Cherchons par exemple $e_3$ tel que $A(e_3) = e_2$ donc en notant $e_3 = \begin{pmatrix}
        x \\ y \\ z
    \end{pmatrix}$, on a le système
    $\begin{cases}
        3x + 9y + 3z = 2 \\
        -3x + 6y + 12z = -1 \\
        x + 8y + 6z = 1
    \end{cases}$ donc $\begin{cases}
        3x + 9y + 3z = 2 \\
        15 y + 15 z = 1 
        30 y + 30 z = 2
    \end{cases}$. 
    
    Choisisson $z = 0$, alors $y = \frac{1}{15}$ et $x = \frac{7}{15}$

    Posons $e_3 = \begin{pmatrix}
        \frac{7}{15} \\ \frac{1}{15} \\ 0
    \end{pmatrix}$.

    Et alors $A e_3 = e_2$.

    $B = (e_1, e_2, e_3)$ est libre et $M_B(u) = \begin{pmatrix}
        15 & 0 & 0 \\
        0 & 0 & 1 \\
        0 & 0 & 0
    \end{pmatrix}$. On a donc bien une matrice triangulaire par blocs. 
\end{exemple}

\subsubsection{Complément HP : décomposition de Dunford }


Soit $u \in \mathscr L(E)$ tel que $\chi_u$ est scindé, 

$\chi_u = \displaystyle \prod_{i = 1 }^{ p } (X - \lambda_i)^{m_i}$ et 
$\forall i \in \ninter{1}{p}, \, C_{\lambda_i}(u) = \ker (u - \lambda_i \id)^{m_i}$.

Alors $E = \displaystyle \bigoplus_{i = 1 }^{ p } C_{\lambda_i}(u)$ et
en notant $\pi_1, \ldots, \pi_p$ les projecteurs associés, il existe un base adaptée $B = B_1 \cup \cdots \cup B_p$
telle que 
$$M_B(u) = \begin{pmatrix}
    M_1 & 0 & \cdots & 0 \\
    0 & M_2 & \cdots & 0 \\
    \vdots & \vdots & \ddots & \vdots \\
    0 & 0 & \cdots & M_p
\end{pmatrix}$$ avec $M_i = \begin{pmatrix}
    \lambda_i & * & \cdots & * \\
    0 & \lambda_i & \cdots & * \\
    \vdots & \vdots & \ddots & \vdots \\
    0 & 0 & \cdots & \lambda_i
\end{pmatrix}$

Notons $d \in \mathscr L(E)$ tel que $M_B(d) = \begin{pmatrix}
    \lambda_1 & 0 & \cdots & \cdots & \cdots & \cdots & 0 \\
    0 & \ddots & \ddots & & & & \vdots \\
    \vdots & \ddots & \lambda_1 & \ddots & & & \vdots \\
    \vdots & & \ddots & \ddots & \ddots &  & \vdots \\
    \vdots &  & & \ddots & \lambda_p &\ddots  & \vdots \\
    \vdots & &  & & & \ddots & 0 \\
    0 & \cdots & \cdots & \cdots & \cdots & 0 & \lambda_p        
\end{pmatrix}$ 

et $n \in \mathscr L(E)$ tel que $n_0 = n - d$. 

Alors $M_B(n_0) = \begin{pmatrix}
    N_1 & 0 & \cdots & 0 \\
    0 & N_2 & \cdots & 0 \\
    \vdots & \vdots & \ddots & \vdots \\
    0 & 0 & \cdots & N_p
\end{pmatrix}$ avec $N_i = \begin{pmatrix}
    0 & * & \cdots & * \\
    0 & 0 & \cdots & * \\
    \vdots & \vdots & \ddots & \vdots \\
    0 & 0 & \cdots & 0
\end{pmatrix}$

alors $\chi_{n_0} = X^n$ donc $n_0$ nilpotent et pour $i \in \ninter 1 p$, 
$\C_{\lambda_i}(u)$ est stable par $n$ et l'endomorphisme induit est 
$n_i$. 

$d$ est diagonalisable, de plus par produit par blocs, 
$n_0 \circ d = d \circ n_0$. 



\begin{theoreme}{Décomposition de Dunford}{}
    Pour $u \in \mathscr L(E)$ tel que $\chi_u$ scindé, alors il existe des endomorphismes 
    $d, n_0 \in \mathscr L(E)$ uniques tels que 
    \begin{itemize}
        \item $d$ diagonalisable
        \item $n_0$ nilpotent
        \item $d \circ n_0 = n_0 \circ d$
        \item $u = d + n_0$
        \item $n_0 \in \K[u]$
        \item $d \in \K[u]$
    \end{itemize}
\end{theoreme}

\begin{proof}
    Existence faite avant 

    Reste donc à démontrer que $n_0$ et $d$ dont des polynômes en $u$.

    Pour cela on montre que $\forall i, $ $\pi_i$ est un polynôme en $u$ 
    (par Bézout). 

    \avspace Unicité : si $d'$ et $n_0'$ conviennent dans $\K[u]$, alors 
    $d + n_0 = d' + n_0'$ donc $d - d' = n_0' - n_0$.

    $d$ et $d' \in \K[u]$ donc comutent. Ils sont diagonalisables donc $d - d'$ aussi.
    $n_0$ et $n_0'$ sont nilpotents donc $n_0' - n_0$ aussi.

    Donc $d - d'$ est diagonalisable et nilpotent donc $d - d' = 0$ et donc $n_0 - n_0' = 0$.
    d'où l'unicité.
\end{proof}

